%% =====================================================================
%%  Introduction to Mathematics for Deep Learning -- Midterm 2026
%%
%%  One source, two outputs.  Choose the version with the switch below,
%%  or from the command line without editing the file:
%%    pdflatex -jobname=midterm26_prob midterm26.tex
%%    pdflatex -jobname=midterm26_sol  "\def\SOLUTION{}%% =====================================================================
%%  Introduction to Mathematics for Deep Learning -- Midterm 2026
%%
%%  One source, two outputs.  Choose the version with the switch below,
%%  or from the command line without editing the file:
%%    pdflatex -jobname=midterm26_prob midterm26.tex
%%    pdflatex -jobname=midterm26_sol  "\def\SOLUTION{}%% =====================================================================
%%  Introduction to Mathematics for Deep Learning -- Midterm 2026
%%
%%  One source, two outputs.  Choose the version with the switch below,
%%  or from the command line without editing the file:
%%    pdflatex -jobname=midterm26_prob midterm26.tex
%%    pdflatex -jobname=midterm26_sol  "\def\SOLUTION{}%% =====================================================================
%%  Introduction to Mathematics for Deep Learning -- Midterm 2026
%%
%%  One source, two outputs.  Choose the version with the switch below,
%%  or from the command line without editing the file:
%%    pdflatex -jobname=midterm26_prob midterm26.tex
%%    pdflatex -jobname=midterm26_sol  "\def\SOLUTION{}\input{midterm26}"
%%  (run twice so "Page x of y" resolves).
%% =====================================================================

\documentclass[12pt]{article}

%% ---------------------------------------------------------------------
%%  0. Version switch: problem sheet vs. solutions
%% ---------------------------------------------------------------------
\newif\ifsolution
\solutionfalse            % \solutiontrue -> solutions version
\ifdefined\SOLUTION \solutiontrue \fi

%% ---------------------------------------------------------------------
%%  1. Page layout (A4)
%% ---------------------------------------------------------------------
\usepackage[a4paper,top=1in,bottom=1in,left=1in,right=1in,footskip=30pt]{geometry}

%% ---------------------------------------------------------------------
%%  2. Packages
%% ---------------------------------------------------------------------
% versions / lists / verbatim
\usepackage{versions}
\usepackage{enumerate}
%\usepackage{enumitem}
\usepackage{verbatim}

% math
\usepackage{amsmath}
\usepackage{amsthm}
\usepackage{amsopn}
\usepackage{amsfonts}

% graphics
\usepackage{graphicx}
\usepackage{tikz}
\usetikzlibrary{arrows,shapes}
\usetikzlibrary{fit}
\usepackage{pgfplots}
\usepackage{psfrag}
\usepackage{pstricks}

% header / footer
\usepackage{fancyhdr}
\usepackage{lastpage}

% references and notes (cleveref must come after hyperref)
\usepackage{hyperref}
\usepackage[capitalize,noabbrev]{cleveref}
% Todonotes is useful during development; to turn off comments use
%\usepackage[disable,textsize=tiny]{todonotes}
\usepackage[textsize=tiny]{todonotes}

\allowdisplaybreaks

%% ---------------------------------------------------------------------
%%  3. Header / footer
%% ---------------------------------------------------------------------
\pagestyle{fancy}
\renewcommand{\headrulewidth}{0pt}
\renewcommand{\footrulewidth}{0.3pt}
\newcommand{\hffont}{\fontfamily{phv}\fontsize{9}{11}\selectfont}
\def\hfsize{\footnotesize}
\def\pagenumber{Page \thepage\ of \pageref{LastPage}}
\fancyhf{}
\newcommand{\LeftFoot}[1]{\fancyfoot[LO]{\textsl{\hffont #1}}}
\fancyfoot[LE]{\textsl{\hffont \pagenumber}}
\fancyfoot[RE]{\textsl{\hffont Introduction to Mathematics for Deep Learning, 2026}}
\fancyfoot[RO]{\textsl{\hffont \pagenumber}}
%\thispagestyle{empty}

%% ---------------------------------------------------------------------
%%  4. Problem / solution environments
%% ---------------------------------------------------------------------
%  \begin{problem}...\end{problem}   : problem sheet only (cover, blank pages)
%  \begin{solution}...\end{solution} : solutions only
\ifsolution
  \excludeversion{problem}
  \includeversion{solution}
\else
  \includeversion{problem}
  \excludeversion{solution}
\fi

\renewenvironment{proof}{\noindent{\bf Proof}\hspace*{1em}}{\qed\bigskip\\}
\newcommand{\sol}{\noindent{\bf Solution: }}

%% ---------------------------------------------------------------------
%%  5. Macros
%% ---------------------------------------------------------------------

% --- 5.1 Calligraphic: \cA ... \cZ, \calD, \calG
\newcommand{\cA}{{\cal A}}
\newcommand{\cB}{{\cal B}}
\newcommand{\cC}{{\cal C}}
\newcommand{\cD}{{\cal D}}
\newcommand{\cE}{{\cal E}}
\newcommand{\cF}{{\cal F}}
\newcommand{\cG}{{\cal G}}
\newcommand{\cH}{{\cal H}}
\newcommand{\cI}{{\cal I}}
\newcommand{\cJ}{{\cal J}}
\newcommand{\cK}{{\cal K}}
\newcommand{\cL}{{\cal L}}
\newcommand{\cM}{{\cal M}}
\newcommand{\cN}{{\cal N}}
\newcommand{\cO}{{\cal O}}
\newcommand{\cP}{{\cal P}}
\newcommand{\cQ}{{\cal Q}}
\newcommand{\cR}{{\cal R}}
\newcommand{\cS}{{\cal S}}
\newcommand{\cT}{{\cal T}}
\newcommand{\cU}{{\cal U}}
\newcommand{\cV}{{\cal V}}
\newcommand{\cW}{{\cal W}}
\newcommand{\cX}{{\cal X}}
\newcommand{\cY}{{\cal Y}}
\newcommand{\cZ}{{\cal Z}}
\newcommand{\mcal}{\mathcal}

\DeclareMathAlphabet{\pazocal}{OMS}{zplm}{m}{n}
\newcommand{\calD}{{\pazocal {D}}}
\newcommand{\calG}{{\pazocal {G}}}

% --- 5.2 Bold lowercase: \ba ... \bz, \eb, \bmu
\newcommand{\ba}{{\bf a}}
\newcommand{\bb}{{\bf b}}
\newcommand{\bc}{{\bf c}}
\newcommand{\bd}{{\bf d}}
\newcommand{\be}{{\bf e}}
\newcommand{\eb}{{\bf e}}
%\newcommand{\bf}{{\bf f}}   % clashes with \bf
\newcommand{\bg}{{\bf g}}
\newcommand{\bh}{{\bf h}}
\newcommand{\bj}{{\bf j}}
\newcommand{\bk}{{\bf k}}
\newcommand{\bl}{{\bf l}}
\newcommand{\bm}{{\bf m}}
\newcommand{\bn}{{\bf n}}
\newcommand{\bo}{{\bf o}}
\newcommand{\bp}{{\bf p}}
\newcommand{\bq}{{\bf q}}
\newcommand{\br}{{\bf r}}
\newcommand{\bs}{{\bf s}}
\newcommand{\bt}{{\bf t}}
\newcommand{\bu}{{\bf u}}
\newcommand{\bv}{{\bf v}}
\newcommand{\bw}{{\bf w}}
\newcommand{\bx}{{\bf x}}
\newcommand{\by}{{\bf y}}
\newcommand{\bz}{{\bf z}}
\newcommand{\bmu}{{\pmb{\mu}}}

% --- 5.3 Bold uppercase: \bA ... \bZ
\newcommand{\bA}{{\bf A}}
\newcommand{\bB}{{\bf B}}
\newcommand{\bC}{{\bf C}}
\newcommand{\bD}{{\bf D}}
\newcommand{\bE}{{\bf E}}
\newcommand{\bF}{{\bf F}}
\newcommand{\bG}{{\bf G}}
\newcommand{\bH}{{\bf H}}
\newcommand{\bI}{{\bf I}}
\newcommand{\bJ}{{\bf J}}
\newcommand{\bK}{{\bf K}}
\newcommand{\bL}{{\bf L}}
\newcommand{\bM}{{\bf M}}
\newcommand{\bN}{{\bf N}}
\newcommand{\bO}{{\bf O}}
\newcommand{\bP}{{\bf P}}
\newcommand{\bQ}{{\bf Q}}
\newcommand{\bR}{{\bf R}}
\newcommand{\bS}{{\bf S}}
\newcommand{\bT}{{\bf T}}
\newcommand{\bU}{{\bf U}}
\newcommand{\bV}{{\bf V}}
\newcommand{\bW}{{\bf W}}
\newcommand{\bX}{{\bf X}}
\newcommand{\bY}{{\bf Y}}
\newcommand{\bZ}{{\bf Z}}

% --- 5.4 Bold vectors with decorations / superscripts
\newcommand{\bXt}{\tilde{\bf X}}
\newcommand{\bXh}{\hat{\bf X}}
\newcommand{\bNsup}[1]{{\bf N}^{(#1)}}
\newcommand{\bNtsup}[1]{\tilde{\bf N}^{(#1)}}
\newcommand{\bZsup}[1]{{\bf Z}^{(#1)}}
\newcommand{\bZtsup}[1]{\tilde{\bf Z}^{(#1)}}
\newcommand{\bXhsup}[1]{\hat{\bf X}^{(#1)}}
\newcommand{\bXtsup}[1]{\tilde{\bf X}^{(#1)}}
\newcommand{\bzerod}{{\bf 0}_d}
\newcommand{\boned}{{\bf 1}_d}
\newcommand{\bone}{{\bf 1}}

% indexed bold / plain symbols X^{(i)} etc.
\newcommand{\Xb}[1]{{\bf X^{(#1)}}}
\newcommand{\Xe}[1]{ X^{(#1)}}
\newcommand{\Bb}[1]{{\bf B^{(#1)}}}
\newcommand{\Be}[1]{ B^{(#1)}}
\newcommand{\Zb}[1]{{\bf Z^{(#1)}}}
\newcommand{\Ze}[1]{ Z^{(#1)}}
\newcommand{\Yb}[1]{{\bf Y^{(#1)}}}
\newcommand{\Ye}[1]{ Y^{(#1)}}
\newcommand{\Ub}[1]{{\bf U^{(#1)}}}
\newcommand{\Ue}[1]{ U^{(#1)}}
\newcommand{\Vb}[1]{{\bf V^{(#1)}}}
\newcommand{\Ve}[1]{ V^{(#1)}}
\newcommand{\Wb}[1]{{\bf W^{(#1)}}}
\newcommand{\We}[1]{ W^{(#1)}}
\newcommand{\Wbtilde}[1]{{\bf \widetilde{W}^{(#1)}}}
\newcommand{\Wetilde}[1]{ \widetilde{W}^{(#1)}}
\newcommand{\Xbh}[1]{{\bf \hat{X}^{(#1)}}}
\newcommand{\mb}[1]{{\bf m^{(#1)}}}
\newcommand{\me}[1]{ m^{(#1)}}
\newcommand{\mbhat}[1]{{\bf \hat{m}^{(#1)}}}
\newcommand{\mehat}[1]{{\hat{m}^{(#1)}}}
\newcommand{\Ab}{{\bf A}}
\newcommand{\Qb}{{\bf Q}}
\newcommand{\Xbhat}{{\bf \hat{X}}}

% --- 5.5 Tilde / hat / star
\newcommand{\tA}{\tilde{A}}
\newcommand{\tB}{\tilde{B}}
\newcommand{\tC}{\tilde{C}}
\newcommand{\tD}{\tilde{D}}
\newcommand{\tE}{\tilde{E}}
\newcommand{\tF}{\tilde{F}}
\newcommand{\tG}{\tilde{G}}
\newcommand{\tH}{\tilde{H}}
\newcommand{\tI}{\tilde{I}}
\newcommand{\tJ}{\tilde{J}}
\newcommand{\tK}{\tilde{K}}
\newcommand{\tL}{\tilde{L}}
\newcommand{\tM}{\tilde{M}}
\newcommand{\tN}{\tilde{N}}
\newcommand{\tO}{\tilde{O}}
\newcommand{\tP}{\tilde{P}}
\newcommand{\tQ}{\tilde{Q}}
\newcommand{\tR}{\tilde{R}}
\newcommand{\tS}{\tilde{S}}
\newcommand{\tT}{\tilde{T}}
\newcommand{\tU}{\tilde{U}}
\newcommand{\tV}{\tilde{V}}
\newcommand{\tX}{\tilde{X}}
\newcommand{\tY}{\tilde{Y}}
\newcommand{\tZ}{\tilde{Z}}
\newcommand{\tu}{\tilde{u}}
\newcommand{\tx}{\tilde{x}}
\newcommand{\ty}{\tilde{y}}
\newcommand{\tgamma}{\tilde{\gamma}}
%\newcommand{\ta}{\tilde{a}}

\newcommand{\uh}{\hat{u}}
\newcommand{\vh}{\hat{v}}
\newcommand{\xh}{\hat{x}}
\newcommand{\yh}{\hat{y}}
\newcommand{\zh}{\hat{z}}
\newcommand{\Ah}{\hat{A}}
\newcommand{\Uh}{\hat{U}}
\newcommand{\Xh}{\hat{X}}
\newcommand{\Yh}{\hat{Y}}
\newcommand{\Zh}{\hat{Z}}
\newcommand{\cXh}{\hat{\cX}}
\newcommand{\cYh}{\hat{\cY}}
\newcommand{\cZh}{\hat{\cZ}}
\newcommand{\bUh}{\hat{\bU}}

\newcommand{\bQstar}{\bQ^{\star}}
\newcommand{\Astar}{A^{\star}}
\newcommand{\Bstar}{B^{\star}}
\newcommand{\Qstar}{Q^{\star}}

% --- 5.6 Blackboard bold (number sets)
\newcommand{\field}[1]{\mathbb{#1}}
\newcommand{\fA}{\field{A}}
\newcommand{\fB}{\field{B}}
\newcommand{\fC}{\field{C}}
\newcommand{\fE}{\field{E}}
\newcommand{\fI}{\field{I}}
\newcommand{\fN}{\field{N}} % natural numbers
\newcommand{\fR}{\field{R}} % real numbers
\newcommand{\fS}{\field{S}}
\newcommand{\fT}{\field{T}}
\newcommand{\fZ}{\field{Z}} % integers

% --- 5.7 Probability: distributions, expectation, relations
\newcommand{\Bern}{{\mbox{Bern}}}
\newcommand{\Lap}{{\mbox{Lap}}}
\newcommand{\SpLap}{{\mbox{SpLap}}}
\newcommand{\punif}{P_{\text{unif}}}
\newcommand{\mathUO}{{\rm Unif}(\mathcal{O}(n))}
\newcommand{\iid}{\stackrel{iid}{\sim}}
\newcommand{\diid}{\overset{\text{i.i.d.}}{\sim}}
\newcommand{\deq}{\stackrel{(d)}{=}}

\newcommand{\mathify}[1]{\ifmmode{#1}\else\mbox{$#1$}\fi}
\newcommand{\Expec}[1]{\mathify{\mathbb{E}\left[#1\right]}}
\newcommand{\Esub}[2]{\mathify{\mathbb{E}_{#1}\left[#2\right]}}
% \newcommand{\Esubs}[2]{\mathify{\mathbb{E}_{#1}\left[#2\right]}}
\newcommand{\Var}[1]{\mathify{{\rm VAR}\left[#1\right]}}
\newcommand{\Cov}{\mathify{{\rm Cov}}}
\newcommand{\Err}{\mathrm{Err}}
\newcommand{\Espace}{{\tilde{\mathcal{F}}[0,T]}}

% --- 5.8 Estimation / text-style operators
\newcommand{\tsnr}{{\text{snr}}}
\newcommand{\tvar}{{\text{Var}}}
\newcommand{\cmle}{\text{\rm cmle}}
\newcommand{\cmse}{\text{\rm cmse}}
\newcommand{\mmse}{\text{\rm mmse}}
\newcommand{\mmle}{\text{\rm mmle}}
\newcommand{\minimax}{\text{\rm minimax}}
\newcommand{\mse}{\text{\rm mse}}
\newcommand{\mle}{\text{\rm mle}}
\newcommand{\diag}{\mbox{diag}}
\newcommand{\JSD}{\mbox{JSD}}
\newcommand{\conv}{\text{\rm conv}}
\newcommand{\sgn}{{\mbox{sign}}}

\newcommand\argmin{\mathop{\mbox{{\rm argmin}}}\limits}
\newcommand\argmax{\mathop{\mbox{{\rm argmax}}}\limits}
\newcommand\maximize{\mathop{\mbox{{\rm maximize}}}\limits}
\newcommand\minimize{\mathop{\mbox{{\rm minimize}}}\limits}

% --- 5.9 Delimiters and relations
\newcommand{\lp}{\left(}
\newcommand{\rp}{\right)}
\newcommand{\lb}{\left[}
\newcommand{\rb}{\right]}
\newcommand{\lbp}{\left\{}
\newcommand{\rbp}{\right\}}
\newcommand{\lba}{\left\lvert}
\newcommand{\rba}{\right\rvert}
\newcommand{\lV}{\left\lVert}
\newcommand{\rV}{\right\rVert}
\newcommand{\mv}{\middle\vert}
\newcommand{\lan}{\langle}
\newcommand{\ran}{\rangle}
\newcommand{\ra}{\rightarrow}
\newcommand{\eqDef}{\coloneqq}   % needs mathtools
\newcommand{\suchthat}{\;\ifnum\currentgrouptype=16 \middle\fi|\;}
\newcommand{\dbar}{\;\ifnum\currentgrouptype=16 \middle \fi\|\;}
\newcommand{\equalto}[2]{\underset{\scriptstyle\overset{\mkern4mu\verteq}{#2}}{#1}}
\newcommand{\verteq}{\rotatebox{90}{$\,=$}}

% --- 5.10 Equation-array shortcuts
\newcommand{\bea}{\begin{eqnarray}}
\newcommand{\eea}{\end{eqnarray}}
\newcommand{\beas}{\begin{eqnarray*}}
\newcommand{\eeas}{\end{eqnarray*}}

% --- 5.11 Comment notes
\newcommand{\noa}[1]{{\color{red}\textbf{AN}: #1}}
\newcommand{\berivan}[1]{{\color{orange}\textbf{BI}: #1}}
\newcommand{\wn}[1]{{\color{green}\textbf{WN}: #1}}

% --- 5.12 Compatibility
% make hyperref and algorithmic work together better
\newcommand{\theHalgorithm}{\arabic{algorithm}}


%% =====================================================================
\begin{document}
%% =====================================================================

%% ---------------------------------------------------------------------
%%  Cover
%% ---------------------------------------------------------------------
\begin{center}
{\bf Introduction to Mathematics for Deep Learning: Midterm \begin{solution}Solutions\end{solution}} \\ \medskip
\begin{problem}Tuesday, October 20, 2026\end{problem}
\end{center}

\begin{problem}
\noindent
\begin{itemize}
\item You can use previous parts of a problem even if you did not solve them.
\item As throughout the course, entropy ($H$) and Mutual Information ($I$) are in bits.
\item $\log$ is taken in base 2; $\ln$ is the natural logarithm.
\end{itemize}
\LeftFoot{Midterm}
\end{problem}

\begin{solution}
\LeftFoot{Midterm Solution}
\end{solution}

\begin{problem}
    \vskip .3in
    {\bf ID : \underline{\hspace{6cm}}}\\
    \vskip .3in
    {\bf Name : \underline{\hspace{5.5cm}}}\\
    ~
    \vskip 2in
    \begin{table}[h]
        \Large
        \begin{tabular}{c|c}
           1     & \textcolor{white}{aaaaaaaa}/~40 \\
           2     & \textcolor{white}{aaaaaaaa}/~30 \\
           3     & \textcolor{white}{aaaaaaaa}/~30 \\
           \hline
           Total & \textcolor{white}{aaaaaaaa}/~100 \\
        \end{tabular}
    \end{table}
    \newpage
    {\bf extra sheet}
    \newpage
\end{problem}

\begin{enumerate}

%% =====================================================================
%%  Problem 1
%% =====================================================================
\item {\bf Mix of Questions} {\it (40 points, 10 points each)}

\begin{enumerate}[a)]

% (a)
\item Let $X_1,\dots,X_n$ be i.i.d.\ $\Bern(\frac12)$ random variables, $X^n=(X_1,\dots,X_n)$, and let $K=\sum_{i=1}^n X_i$ be the number of ones. Show that
\[
H(X^n\mid K)\ \ge\ n-\log(n+1).
\]

% (b)
\item Let $X$ and $Y$ be continuous random variables with densities $f_X$ and $f_Y$, respectively, and let $Z$ have density $f_Z=\frac12 f_X+\frac12 f_Y$. Prove or disprove: $h(Z)\le\max\{h(X),\,h(Y)\}$.

% (c)
\item Let $X$ be uniform on $K\ge 2$ classes $\{1,\dots,K\}$. A noisy label $Y$ is generated as follows: with probability $1-\varepsilon$, $Y=X$; with probability $\varepsilon$, $Y$ is drawn uniformly from the remaining $K-1$ classes (all classes other than $X$).
\begin{enumerate}[i)]
\item Compute $I(X;Y)$ for $\varepsilon=0$ and for $\varepsilon=1$, and show that both are nonzero.
\item Find $\varepsilon$ such that $I(X;Y)=0$.
\end{enumerate}

% (d)
\item A deck of $n$ cards is shuffled by the following move: take the top card and reinsert it at one of the $n$ possible positions (on top, below the 1st card, below the 2nd card, $\dots$, at the bottom), chosen uniformly at random and independently of all previous moves. Reinserting at the top means the card is put back where it was, so the deck stays the same. Let $X_0$ be the initial ordering of the deck and $X_i$ the ordering of the $n$ cards after the $i$-th move. Then $X_0,X_1,X_2,\dots$ is a Markov chain on the $n!$ orderings of the deck. Show that the uniform distribution over all $n!$ orderings is a stationary distribution of this Markov chain.
\end{enumerate}

%%% Solution of Problem 1
\begin{solution}
\vskip 0.2in
{\bf Solution:}

\begin{enumerate}[a)]

%% (a) sol
\item Since $K$ is a function of $X^n$, $H(X^n,K)=H(X^n)$. By the chain rule,
\[
n=H(X^n)=H(X^n,K)=H(K)+H(X^n\mid K).
\]
$K$ takes $n+1$ values $\{0,1,\dots,n\}$, so $H(K)\le\log(n+1)$. Hence
\[
H(X^n\mid K)=n-H(K)\ \ge\ n-\log(n+1).
\]
(Remark: the count $K$ carries at most $\log(n+1)$ of the $n$ bits; almost all of the information is in the order of the bits.)

%% (b) sol
\item The statement is {\bf false}. Counterexample: $X\sim{\rm Unif}[0,1]$, $Y\sim{\rm Unif}[10,11]$. Then
\[
h(X)=-\int_0^1 1\cdot\log 1\,dx=0,\qquad h(Y)=0.
\]
$Z$ has density $f_Z(z)=\frac12$ on $[0,1]\cup[10,11]$ (total length 2) and $0$ elsewhere, so
\[
h(Z)=-\int f_Z\log f_Z=-2\cdot\tfrac12\log\tfrac12=1\ >\ 0=\max\{h(X),h(Y)\}.
\]
(Remark: mixing adds the uncertainty about \emph{which} density $Z$ came from. Introducing a label $W\sim\Bern(\frac12)$ with $Z\mid\{W=1\}\sim f_X$ and $Z\mid\{W=0\}\sim f_Y$, one can show that, whenever $h(X)$ and $h(Y)$ are finite,
\[
\tfrac12h(X)+\tfrac12h(Y)\ \le\ h(Z)\ \le\ \tfrac12h(X)+\tfrac12h(Y)+1,
\]
with the upper bound attained exactly when $f_X(z)f_Y(z)=0$ for almost every $z$, as above.)

%% (c) sol
\item The channel is $P(Y=y\mid X=x)=1-\varepsilon$ if $y=x$ and $\frac{\varepsilon}{K-1}$ if $y\ne x$. By symmetry $Y$ is uniform, so $H(Y)=\log K$ and $H(X)=\log K$.
\begin{enumerate}[i)]
\item $\varepsilon=0$: $Y=X$, so $I(X;Y)=H(X)-H(X\mid Y)=\log K-0=\log K>0$.

$\varepsilon=1$: $Y$ is always wrong. Given $Y=y$, $X$ is uniform over the $K-1$ classes other than $y$, so $H(X\mid Y)=\log(K-1)$ and
\[
I(X;Y)=\log K-\log(K-1)=\log\frac{K}{K-1}>0.
\]
A label that is always wrong still tells us which class $X$ is \emph{not}. (For $K=2$ this is a full bit, $Y=3-X$; for $K\ge 3$ it is less than $\log K$.) Mutual information measures dependence, not correctness.

\item $I(X;Y)=0$ iff $X$ and $Y$ are independent, i.e., iff the conditional law $P(Y=\cdot\mid X=x)$ is the same for every $x$. That law puts $1-\varepsilon$ on $x$ and $\frac{\varepsilon}{K-1}$ on every other class, so it does not depend on $x$ iff it is uniform:
\[
1-\varepsilon=\frac{\varepsilon}{K-1}\iff\varepsilon=\frac{K-1}{K}.
\]
Then $Y$ is uniform regardless of $X$, so $I(X;Y)=0$. (Note: this is $\varepsilon=\frac12$ only for $K=2$; for $K=10$, labels that are wrong $60\%$ of the time are still informative.)
\end{enumerate}

%% (d) sol
\item Let $P(\tau\to\sigma)=P(X_{i+1}=\sigma\mid X_i=\tau)$ be the one-step transition probability between orderings and $U(\sigma)=\frac1{n!}$. We need $\sum_\tau U(\tau)P(\tau\to\sigma)=U(\sigma)$ for every ordering $\sigma$.
\begin{itemize}
\item From any ordering $\tau$, the $n$ insertion positions of the top card give $n$ \emph{different} orderings, each with probability $\frac1n$. So $P(\tau\to\sigma)\in\{0,\frac1n\}$.
\item Which $\tau$ can reach a fixed $\sigma$? Undo the move: pick any one of the $n$ cards of $\sigma$ and move it to the top. This gives exactly $n$ distinct predecessors $\tau$ (picking the top card of $\sigma$ gives $\tau=\sigma$, the ``reinsert on top'' move).
\end{itemize}
Therefore
\[
\sum_\tau U(\tau)P(\tau\to\sigma)=n\cdot\frac1{n!}\cdot\frac1n=\frac1{n!}=U(\sigma).
\]
(Equivalently, the transition matrix is doubly stochastic: every row and every column sums to 1, so the uniform distribution is stationary.)
\end{enumerate}
\end{solution}

\begin{problem}
\newpage {\bf problem 1 cont.}
\end{problem}
\newpage


%% =====================================================================
%%  Problem 2
%% =====================================================================
\item {\bf Label Smoothing vs.\ Confidence Penalty} {\it (30 points)}

A classifier for $K\ge 2$ classes outputs a probability distribution $q=(q_1,\dots,q_K)$ with all $q_k>0$, where $q_k$ is the probability it assigns to label $k$. Let $y$ be the true label, and let $u=(\frac1K,\dots,\frac1K)$ be the uniform distribution. Two common ways to discourage over-confident predictions are:
\begin{itemize}
\item {\bf Label smoothing.} For $0<\varepsilon<1$, replace the one-hot target by the \emph{smoothed one-hot} vector $o^{(\varepsilon)}=(o^{(\varepsilon)}_1,\dots,o^{(\varepsilon)}_K)$ with
\[
o^{(\varepsilon)}_k=
\begin{cases}
1-\varepsilon+\dfrac{\varepsilon}{K}, & k=y \ \text{(the true label)},\\[2mm]
\dfrac{\varepsilon}{K}, & k\ne y,
\end{cases}
\]
and minimize the cross-entropy $L_{\rm LS}(q)=-\sum_{k=1}^K o^{(\varepsilon)}_k\log q_k$. Note that $o^{(0)}$ is the one-hot vector of the true label $y$ (a $1$ in position $y$ and $0$ elsewhere), so $\varepsilon=0$ gives the usual cross-entropy loss $-\log q_y$.
\item {\bf Confidence penalty.} For $\beta>0$, minimize $L_{\rm CP}(q)=-\log q_y-\beta H(q)$, where $H(q)=-\sum_k q_k\log q_k$.
\end{itemize}

\begin{enumerate}[(a)]
\item Show the following: {\it (10 points)}
\begin{align*}
L_{\rm LS}(q)&=(1-\varepsilon)(-\log q_y)+\varepsilon\,D(u\|q)+\varepsilon\log K,\\
L_{\rm CP}(q)&=-\log q_y+\beta\,D(q\|u)-\beta\log K.
\end{align*}

\item Fix the true label $y$ and let the prediction become fully confident: $q_y\to 1$ (and of course $q_k\to 0$ for every $k\ne y$). Find the limits of $D(u\|q)$, $D(q\|u)$, $L_{\rm LS}(q)$ and $L_{\rm CP}(q)$. {\it (10 points)}

\item In this part only, $q$ may have zero entries. Let $p$ be any distribution on $\{1,\dots,K\}$.
\begin{enumerate}[i)]
\item Show that $D(p\|q)<\infty$ if and only if $q_k>0$ for every $k$ with $p_k>0$. {\it (4 points)}
\item Show that $D(q\|p)<\infty$ if and only if $q_k=0$ for every $k$ with $p_k=0$. {\it (4 points)}
\item Set $p=u$. For which $q$ is $D(u\|q)$ finite? For which $q$ is $D(q\|u)$ finite? {\it (2 points)}
\end{enumerate}
\end{enumerate}

%%% Solution of Problem 2
\begin{solution}
\vskip 0.2in
{\bf Solution:}

\begin{enumerate}[(a)]

%% (a) sol
\item Write $o^{(\varepsilon)}=(1-\varepsilon)o^{(0)}+\varepsilon u$, where $o^{(0)}$ is the one-hot vector of $y$. The cross-entropy is linear in the target, so
\[
L_{\rm LS}(q)=(1-\varepsilon)(-\log q_y)+\varepsilon\Big(-\sum_k\tfrac1K\log q_k\Big).
\]
For the second term,
\[
-\sum_k\tfrac1K\log q_k=\sum_k\tfrac1K\log\frac{1/K}{q_k}-\sum_k\tfrac1K\log\tfrac1K=D(u\|q)+\log K.
\]
This gives the first identity. For the second,
\[
D(q\|u)=\sum_k q_k\log\frac{q_k}{1/K}=\sum_k q_k\log q_k+\log K=\log K-H(q),
\]
so $-\beta H(q)=\beta D(q\|u)-\beta\log K$, which gives the second identity.

%% (b) sol
\item $D(u\|q)=\frac1K\sum_k\log\frac{1}{Kq_k}$. Each term with $k\ne y$ tends to $+\infty$ as $q_k\to 0$, so $D(u\|q)\to\infty$.

$D(q\|u)=\log K-H(q)\to\log K-0=\log K$, which is finite.

Hence $L_{\rm LS}(q)\to\infty$, since $(1-\varepsilon)(-\log q_y)\to 0$ but $\varepsilon D(u\|q)\to\infty$; and $L_{\rm CP}(q)\to 0+\beta\log K-\beta\log K=0$. Full confidence is infinitely expensive under label smoothing, but has finite loss under the confidence penalty.

%% (c) sol
\item Use the conventions $0\log\frac{0}{a}=0$ for $a\ge 0$ and $b\log\frac{b}{0}=+\infty$ for $b>0$.
\begin{enumerate}[i)]
\item $D(p\|q)=\sum_k p_k\log\frac{p_k}{q_k}$ is a finite sum. Terms with $p_k=0$ contribute $0$; a term with $p_k>0$ is finite if $q_k>0$ and $+\infty$ if $q_k=0$. Hence $D(p\|q)<\infty$ iff {\bf $q_k>0$ for every $k$ with $p_k>0$}.
\item $D(q\|p)=\sum_k q_k\log\frac{q_k}{p_k}$. Terms with $q_k=0$ contribute $0$; a term with $q_k>0$ is finite if $p_k>0$ and $+\infty$ if $p_k=0$. Hence $D(q\|p)<\infty$ iff {\bf $q_k=0$ for every $k$ with $p_k=0$}.
\item Since $u_k=\frac1K>0$ for every $k$:
\begin{itemize}
\item by (i), $D(u\|q)<\infty$ iff {\bf $q_k>0$ for every $k$} (and $D(u\|q)\to\infty$ as any $q_k\to 0$);
\item by (ii), the condition is empty, so $D(q\|u)<\infty$ for {\bf every} $q$; indeed $0\le D(q\|u)=\log K-H(q)\le\log K$.
\end{itemize}
\end{enumerate}
(Remark: this explains (b). $L_{\rm LS}$ contains $\varepsilon D(u\|q)$, which puts an unbounded cost on any class probability approaching $0$, so $L_{\rm LS}\to\infty$ as $q_y\to 1$. $L_{\rm CP}$ contains $\beta D(q\|u)$, which is bounded, so the fully confident limit has finite loss. Neither is a hard constraint: e.g.\ $L_{\rm LS}$ is minimized at $q=o^{(\varepsilon)}$, but any $q$ with all $q_k>0$ has finite loss. Minimizing $D(p\|q)$ over $q$ is called \emph{zero-avoiding} and minimizing $D(q\|p)$ \emph{zero-forcing}, after (i) and (ii).)
\end{enumerate}
\end{solution}

\begin{problem}
\newpage {\bf problem 2 cont.}
\end{problem}
\newpage


%% =====================================================================
%%  Problem 3
%% =====================================================================
\item {\bf Heavy Tails and the Median of Means} {\it (30 points)}

Let $X$ be a random variable with $X\ge 1$ and
\[
P(X>x)=x^{-3}\qquad\text{for all } x\ge 1.
\]
Let $X_1,\dots,X_n$ be i.i.d.\ copies of $X$, and $\bar X_n=\frac1n\sum_{i=1}^n X_i$. One can check that $\mathbb{E}[e^{tX}]=\infty$ for every $t>0$, so the Chernoff bound cannot be applied directly to $\bar X_n$.

{\it Hint.} For a nonnegative random variable $W$, $\mathbb{E}[W]=\int_0^\infty P(W>w)\,dw$. You may use this without proof.

\begin{enumerate}[(a)]
\item Show that $\mathbb{E}[X^2]=\int_0^\infty 2x\,P(X>x)\,dx$. Then compute $\mu=\mathbb{E}[X]$ and $\sigma^2={\rm Var}(X)$. {\it (5 points)}

\item For $a>0$, show that
\[
P(\bar X_n\ge\mu+a)\ \ge\ \big(n(\mu+a)\big)^{-3},
\]
and compare it with the upper bound given by Chebyshev's inequality. Conclude that $P(\bar X_n\ge\mu+a)$ decays polynomially, not exponentially, in $n$. [Hint: $\bar X_n\ge X_1/n$.] {\it (5 points)}

\item \emph{Median of means.} Split the $n$ samples into $k$ groups of $m$ samples each ($n=km$), compute the sample mean $\bar X^{(j)}$ of each group $j=1,\dots,k$, and let $\hat\mu$ be the median of the $k$ group means. Fix $a>0$ and call group $j$ \emph{bad} if $|\bar X^{(j)}-\mu|\ge a$. {\it (15 points)}
\begin{enumerate}[i)]
\item Show that if $m\ge 6/a^2$, each group is bad with probability at most $\frac18$. [Hint: Chebyshev's inequality for $\bar X^{(j)}$.] {\it (3 points)}
\item Show that if $|\hat\mu-\mu|\ge a$, then at least $k/2$ groups are bad. [Hint: if fewer than half of the groups are bad, where do more than half of the group means lie, and where is their median?] {\it (4 points)}
\item Let $Z$ be the number of bad groups. Since the groups use disjoint samples, you may assume the worst case $Z\sim{\rm Binomial}(k,\frac18)$. Apply the Chernoff bound to show that for every $t>0$,
\[
P(Z\ge k/2)\ \le\ \Big(e^{-t/2}\big(\tfrac78+\tfrac18 e^{t}\big)\Big)^k.
\]
[Hint: what is the moment generating function of a Bernoulli random variable?] {\it (4 points)}
\item Choose $e^t=7$ and show that the bound becomes $\big(\frac{\sqrt7}{4}\big)^k$. Combine (ii)--(iv) to conclude that with $m=\lceil 6/a^2\rceil$, $P(|\hat\mu-\mu|\ge a)$ decays exponentially in $n$. {\it (4 points)}
\end{enumerate}

\item Explain in two or three sentences why the median of means succeeds where the sample mean $\bar X_n$ fails. {\it (5 points)}
\end{enumerate}

%%% Solution of Problem 3
\begin{solution}
\vskip 0.2in
{\bf Solution:}

\begin{enumerate}[(a)]

%% (a) sol
\item Apply the hint to $W=X^2\ge 0$ and substitute $s=x^2$ ($ds=2x\,dx$):
\[
\mathbb{E}[X^2]=\int_0^\infty P(X^2>s)\,ds=\int_0^\infty P(X>\sqrt s)\,ds=\int_0^\infty 2x\,P(X>x)\,dx.
\]
Since $P(X>x)=1$ for $x<1$,
\[
\mu=1+\int_1^\infty x^{-3}\,dx=1+\tfrac12=\tfrac32,\qquad
\mathbb{E}[X^2]=1+\int_1^\infty 2x^{-2}\,dx=1+2=3,
\]
so $\sigma^2=3-\frac94=\frac34$. (Equivalently, use the density $f(x)=3x^{-4}$ for $x\ge 1$.)

%% (b) sol
\item All $X_i\ge 1>0$, so $\bar X_n\ge X_1/n$. If $X_1\ge n(\mu+a)$, then $\bar X_n\ge\mu+a$. Hence
\[
P(\bar X_n\ge\mu+a)\ \ge\ P\big(X_1\ge n(\mu+a)\big)=\big(n(\mu+a)\big)^{-3}.
\]
Chebyshev gives the upper bound
\[
P(\bar X_n\ge\mu+a)\le P(|\bar X_n-\mu|\ge a)\le\frac{\sigma^2}{na^2}=\frac{3}{4na^2}.
\]
So for fixed $a>0$ the probability tends to $0$ but is bounded above and below by polynomial rates ($n^{-1}$ and $n^{-3}$). In particular no bound $Ce^{-cn}$ with fixed $C,c>0$ can hold for all large $n$, because $(n(\mu+a))^{-3}\gg Ce^{-cn}$.

%% (c) sol
\item
\begin{enumerate}[i)]
\item $\bar X^{(j)}$ is the mean of $m$ i.i.d.\ copies of $X$, so ${\rm Var}(\bar X^{(j)})=\sigma^2/m$. By Chebyshev,
\[
P(|\bar X^{(j)}-\mu|\ge a)\le\frac{\sigma^2}{ma^2}=\frac{3}{4ma^2}\le\frac34\cdot\frac16=\frac18\qquad\text{when } m\ge 6/a^2.
\]
\item Suppose fewer than $k/2$ groups are bad. Then more than half of the group means lie in the interval $(\mu-a,\mu+a)$. If $\hat\mu\ge\mu+a$, then more than half of the group means would lie strictly below $\hat\mu$, contradicting the fact that a median has at least half of the values at or above it. The case $\hat\mu\le\mu-a$ is symmetric. Hence $|\hat\mu-\mu|<a$. Contrapositive: $|\hat\mu-\mu|\ge a$ implies at least $k/2$ bad groups, i.e., $\{|\hat\mu-\mu|\ge a\}\subseteq\{Z\ge k/2\}$.
\item Chernoff bound: for $t>0$, $e^{tz}$ is increasing, so by Markov's inequality
\[
P(Z\ge k/2)=P(e^{tZ}\ge e^{tk/2})\le e^{-tk/2}\,\mathbb{E}[e^{tZ}].
\]
$Z=\sum_{j=1}^k B_j$ with independent $B_j\sim\Bern(\frac18)$, and the MGF of a $\Bern(\frac18)$ variable is $\mathbb{E}[e^{tB_j}]=\frac78+\frac18e^t$. By independence, $\mathbb{E}[e^{tZ}]=\prod_j\mathbb{E}[e^{tB_j}]=\big(\frac78+\frac18e^t\big)^k$. Combining,
\[
P(Z\ge k/2)\le\Big(e^{-t/2}\big(\tfrac78+\tfrac18e^t\big)\Big)^k.
\]
\item With $e^t=7$: $e^{-t/2}\big(\frac78+\frac78\big)=\frac1{\sqrt7}\cdot\frac74=\frac{\sqrt7}{4}\approx 0.661<1$. By (ii) and (iii),
\[
P(|\hat\mu-\mu|\ge a)\le P(Z\ge k/2)\le\Big(\frac{\sqrt7}{4}\Big)^k.
\]
(If the bad-group probability is $q<\frac18$ rather than exactly $\frac18$, $Z$ is only smaller, so the bound still holds.) With $m=\lceil 6/a^2\rceil$ fixed, $k=n/m$ grows linearly in $n$, so the bound $\big(\frac{\sqrt7}{4}\big)^{n/m}$ decays exponentially in $n$. (The choice $e^t=7$ is the one that minimizes the bound in (iii).)
\end{enumerate}

%% (d) sol
\item A single sufficiently large sample can push $\bar X_n$ above $\mu+a$, and with a heavy tail such a sample appears with only polynomially small probability. The median of means confines any huge sample to a single group, where it spoils only one vote. The median fails only if half of the independent groups fail at once, which is a coin-flip large deviation with an exponentially small probability.
\end{enumerate}
\end{solution}

\begin{problem}
\newpage {\bf problem 3 cont.}
\end{problem}
\begin{problem}
\newpage {\bf extra sheet}
\newpage {\bf extra sheet}
\end{problem}

\end{enumerate}

\end{document}
"
%%  (run twice so "Page x of y" resolves).
%% =====================================================================

\documentclass[12pt]{article}

%% ---------------------------------------------------------------------
%%  0. Version switch: problem sheet vs. solutions
%% ---------------------------------------------------------------------
\newif\ifsolution
\solutionfalse            % \solutiontrue -> solutions version
\ifdefined\SOLUTION \solutiontrue \fi

%% ---------------------------------------------------------------------
%%  1. Page layout (A4)
%% ---------------------------------------------------------------------
\usepackage[a4paper,top=1in,bottom=1in,left=1in,right=1in,footskip=30pt]{geometry}

%% ---------------------------------------------------------------------
%%  2. Packages
%% ---------------------------------------------------------------------
% versions / lists / verbatim
\usepackage{versions}
\usepackage{enumerate}
%\usepackage{enumitem}
\usepackage{verbatim}

% math
\usepackage{amsmath}
\usepackage{amsthm}
\usepackage{amsopn}
\usepackage{amsfonts}

% graphics
\usepackage{graphicx}
\usepackage{tikz}
\usetikzlibrary{arrows,shapes}
\usetikzlibrary{fit}
\usepackage{pgfplots}
\usepackage{psfrag}
\usepackage{pstricks}

% header / footer
\usepackage{fancyhdr}
\usepackage{lastpage}

% references and notes (cleveref must come after hyperref)
\usepackage{hyperref}
\usepackage[capitalize,noabbrev]{cleveref}
% Todonotes is useful during development; to turn off comments use
%\usepackage[disable,textsize=tiny]{todonotes}
\usepackage[textsize=tiny]{todonotes}

\allowdisplaybreaks

%% ---------------------------------------------------------------------
%%  3. Header / footer
%% ---------------------------------------------------------------------
\pagestyle{fancy}
\renewcommand{\headrulewidth}{0pt}
\renewcommand{\footrulewidth}{0.3pt}
\newcommand{\hffont}{\fontfamily{phv}\fontsize{9}{11}\selectfont}
\def\hfsize{\footnotesize}
\def\pagenumber{Page \thepage\ of \pageref{LastPage}}
\fancyhf{}
\newcommand{\LeftFoot}[1]{\fancyfoot[LO]{\textsl{\hffont #1}}}
\fancyfoot[LE]{\textsl{\hffont \pagenumber}}
\fancyfoot[RE]{\textsl{\hffont Introduction to Mathematics for Deep Learning, 2026}}
\fancyfoot[RO]{\textsl{\hffont \pagenumber}}
%\thispagestyle{empty}

%% ---------------------------------------------------------------------
%%  4. Problem / solution environments
%% ---------------------------------------------------------------------
%  \begin{problem}...\end{problem}   : problem sheet only (cover, blank pages)
%  \begin{solution}...\end{solution} : solutions only
\ifsolution
  \excludeversion{problem}
  \includeversion{solution}
\else
  \includeversion{problem}
  \excludeversion{solution}
\fi

\renewenvironment{proof}{\noindent{\bf Proof}\hspace*{1em}}{\qed\bigskip\\}
\newcommand{\sol}{\noindent{\bf Solution: }}

%% ---------------------------------------------------------------------
%%  5. Macros
%% ---------------------------------------------------------------------

% --- 5.1 Calligraphic: \cA ... \cZ, \calD, \calG
\newcommand{\cA}{{\cal A}}
\newcommand{\cB}{{\cal B}}
\newcommand{\cC}{{\cal C}}
\newcommand{\cD}{{\cal D}}
\newcommand{\cE}{{\cal E}}
\newcommand{\cF}{{\cal F}}
\newcommand{\cG}{{\cal G}}
\newcommand{\cH}{{\cal H}}
\newcommand{\cI}{{\cal I}}
\newcommand{\cJ}{{\cal J}}
\newcommand{\cK}{{\cal K}}
\newcommand{\cL}{{\cal L}}
\newcommand{\cM}{{\cal M}}
\newcommand{\cN}{{\cal N}}
\newcommand{\cO}{{\cal O}}
\newcommand{\cP}{{\cal P}}
\newcommand{\cQ}{{\cal Q}}
\newcommand{\cR}{{\cal R}}
\newcommand{\cS}{{\cal S}}
\newcommand{\cT}{{\cal T}}
\newcommand{\cU}{{\cal U}}
\newcommand{\cV}{{\cal V}}
\newcommand{\cW}{{\cal W}}
\newcommand{\cX}{{\cal X}}
\newcommand{\cY}{{\cal Y}}
\newcommand{\cZ}{{\cal Z}}
\newcommand{\mcal}{\mathcal}

\DeclareMathAlphabet{\pazocal}{OMS}{zplm}{m}{n}
\newcommand{\calD}{{\pazocal {D}}}
\newcommand{\calG}{{\pazocal {G}}}

% --- 5.2 Bold lowercase: \ba ... \bz, \eb, \bmu
\newcommand{\ba}{{\bf a}}
\newcommand{\bb}{{\bf b}}
\newcommand{\bc}{{\bf c}}
\newcommand{\bd}{{\bf d}}
\newcommand{\be}{{\bf e}}
\newcommand{\eb}{{\bf e}}
%\newcommand{\bf}{{\bf f}}   % clashes with \bf
\newcommand{\bg}{{\bf g}}
\newcommand{\bh}{{\bf h}}
\newcommand{\bj}{{\bf j}}
\newcommand{\bk}{{\bf k}}
\newcommand{\bl}{{\bf l}}
\newcommand{\bm}{{\bf m}}
\newcommand{\bn}{{\bf n}}
\newcommand{\bo}{{\bf o}}
\newcommand{\bp}{{\bf p}}
\newcommand{\bq}{{\bf q}}
\newcommand{\br}{{\bf r}}
\newcommand{\bs}{{\bf s}}
\newcommand{\bt}{{\bf t}}
\newcommand{\bu}{{\bf u}}
\newcommand{\bv}{{\bf v}}
\newcommand{\bw}{{\bf w}}
\newcommand{\bx}{{\bf x}}
\newcommand{\by}{{\bf y}}
\newcommand{\bz}{{\bf z}}
\newcommand{\bmu}{{\pmb{\mu}}}

% --- 5.3 Bold uppercase: \bA ... \bZ
\newcommand{\bA}{{\bf A}}
\newcommand{\bB}{{\bf B}}
\newcommand{\bC}{{\bf C}}
\newcommand{\bD}{{\bf D}}
\newcommand{\bE}{{\bf E}}
\newcommand{\bF}{{\bf F}}
\newcommand{\bG}{{\bf G}}
\newcommand{\bH}{{\bf H}}
\newcommand{\bI}{{\bf I}}
\newcommand{\bJ}{{\bf J}}
\newcommand{\bK}{{\bf K}}
\newcommand{\bL}{{\bf L}}
\newcommand{\bM}{{\bf M}}
\newcommand{\bN}{{\bf N}}
\newcommand{\bO}{{\bf O}}
\newcommand{\bP}{{\bf P}}
\newcommand{\bQ}{{\bf Q}}
\newcommand{\bR}{{\bf R}}
\newcommand{\bS}{{\bf S}}
\newcommand{\bT}{{\bf T}}
\newcommand{\bU}{{\bf U}}
\newcommand{\bV}{{\bf V}}
\newcommand{\bW}{{\bf W}}
\newcommand{\bX}{{\bf X}}
\newcommand{\bY}{{\bf Y}}
\newcommand{\bZ}{{\bf Z}}

% --- 5.4 Bold vectors with decorations / superscripts
\newcommand{\bXt}{\tilde{\bf X}}
\newcommand{\bXh}{\hat{\bf X}}
\newcommand{\bNsup}[1]{{\bf N}^{(#1)}}
\newcommand{\bNtsup}[1]{\tilde{\bf N}^{(#1)}}
\newcommand{\bZsup}[1]{{\bf Z}^{(#1)}}
\newcommand{\bZtsup}[1]{\tilde{\bf Z}^{(#1)}}
\newcommand{\bXhsup}[1]{\hat{\bf X}^{(#1)}}
\newcommand{\bXtsup}[1]{\tilde{\bf X}^{(#1)}}
\newcommand{\bzerod}{{\bf 0}_d}
\newcommand{\boned}{{\bf 1}_d}
\newcommand{\bone}{{\bf 1}}

% indexed bold / plain symbols X^{(i)} etc.
\newcommand{\Xb}[1]{{\bf X^{(#1)}}}
\newcommand{\Xe}[1]{ X^{(#1)}}
\newcommand{\Bb}[1]{{\bf B^{(#1)}}}
\newcommand{\Be}[1]{ B^{(#1)}}
\newcommand{\Zb}[1]{{\bf Z^{(#1)}}}
\newcommand{\Ze}[1]{ Z^{(#1)}}
\newcommand{\Yb}[1]{{\bf Y^{(#1)}}}
\newcommand{\Ye}[1]{ Y^{(#1)}}
\newcommand{\Ub}[1]{{\bf U^{(#1)}}}
\newcommand{\Ue}[1]{ U^{(#1)}}
\newcommand{\Vb}[1]{{\bf V^{(#1)}}}
\newcommand{\Ve}[1]{ V^{(#1)}}
\newcommand{\Wb}[1]{{\bf W^{(#1)}}}
\newcommand{\We}[1]{ W^{(#1)}}
\newcommand{\Wbtilde}[1]{{\bf \widetilde{W}^{(#1)}}}
\newcommand{\Wetilde}[1]{ \widetilde{W}^{(#1)}}
\newcommand{\Xbh}[1]{{\bf \hat{X}^{(#1)}}}
\newcommand{\mb}[1]{{\bf m^{(#1)}}}
\newcommand{\me}[1]{ m^{(#1)}}
\newcommand{\mbhat}[1]{{\bf \hat{m}^{(#1)}}}
\newcommand{\mehat}[1]{{\hat{m}^{(#1)}}}
\newcommand{\Ab}{{\bf A}}
\newcommand{\Qb}{{\bf Q}}
\newcommand{\Xbhat}{{\bf \hat{X}}}

% --- 5.5 Tilde / hat / star
\newcommand{\tA}{\tilde{A}}
\newcommand{\tB}{\tilde{B}}
\newcommand{\tC}{\tilde{C}}
\newcommand{\tD}{\tilde{D}}
\newcommand{\tE}{\tilde{E}}
\newcommand{\tF}{\tilde{F}}
\newcommand{\tG}{\tilde{G}}
\newcommand{\tH}{\tilde{H}}
\newcommand{\tI}{\tilde{I}}
\newcommand{\tJ}{\tilde{J}}
\newcommand{\tK}{\tilde{K}}
\newcommand{\tL}{\tilde{L}}
\newcommand{\tM}{\tilde{M}}
\newcommand{\tN}{\tilde{N}}
\newcommand{\tO}{\tilde{O}}
\newcommand{\tP}{\tilde{P}}
\newcommand{\tQ}{\tilde{Q}}
\newcommand{\tR}{\tilde{R}}
\newcommand{\tS}{\tilde{S}}
\newcommand{\tT}{\tilde{T}}
\newcommand{\tU}{\tilde{U}}
\newcommand{\tV}{\tilde{V}}
\newcommand{\tX}{\tilde{X}}
\newcommand{\tY}{\tilde{Y}}
\newcommand{\tZ}{\tilde{Z}}
\newcommand{\tu}{\tilde{u}}
\newcommand{\tx}{\tilde{x}}
\newcommand{\ty}{\tilde{y}}
\newcommand{\tgamma}{\tilde{\gamma}}
%\newcommand{\ta}{\tilde{a}}

\newcommand{\uh}{\hat{u}}
\newcommand{\vh}{\hat{v}}
\newcommand{\xh}{\hat{x}}
\newcommand{\yh}{\hat{y}}
\newcommand{\zh}{\hat{z}}
\newcommand{\Ah}{\hat{A}}
\newcommand{\Uh}{\hat{U}}
\newcommand{\Xh}{\hat{X}}
\newcommand{\Yh}{\hat{Y}}
\newcommand{\Zh}{\hat{Z}}
\newcommand{\cXh}{\hat{\cX}}
\newcommand{\cYh}{\hat{\cY}}
\newcommand{\cZh}{\hat{\cZ}}
\newcommand{\bUh}{\hat{\bU}}

\newcommand{\bQstar}{\bQ^{\star}}
\newcommand{\Astar}{A^{\star}}
\newcommand{\Bstar}{B^{\star}}
\newcommand{\Qstar}{Q^{\star}}

% --- 5.6 Blackboard bold (number sets)
\newcommand{\field}[1]{\mathbb{#1}}
\newcommand{\fA}{\field{A}}
\newcommand{\fB}{\field{B}}
\newcommand{\fC}{\field{C}}
\newcommand{\fE}{\field{E}}
\newcommand{\fI}{\field{I}}
\newcommand{\fN}{\field{N}} % natural numbers
\newcommand{\fR}{\field{R}} % real numbers
\newcommand{\fS}{\field{S}}
\newcommand{\fT}{\field{T}}
\newcommand{\fZ}{\field{Z}} % integers

% --- 5.7 Probability: distributions, expectation, relations
\newcommand{\Bern}{{\mbox{Bern}}}
\newcommand{\Lap}{{\mbox{Lap}}}
\newcommand{\SpLap}{{\mbox{SpLap}}}
\newcommand{\punif}{P_{\text{unif}}}
\newcommand{\mathUO}{{\rm Unif}(\mathcal{O}(n))}
\newcommand{\iid}{\stackrel{iid}{\sim}}
\newcommand{\diid}{\overset{\text{i.i.d.}}{\sim}}
\newcommand{\deq}{\stackrel{(d)}{=}}

\newcommand{\mathify}[1]{\ifmmode{#1}\else\mbox{$#1$}\fi}
\newcommand{\Expec}[1]{\mathify{\mathbb{E}\left[#1\right]}}
\newcommand{\Esub}[2]{\mathify{\mathbb{E}_{#1}\left[#2\right]}}
% \newcommand{\Esubs}[2]{\mathify{\mathbb{E}_{#1}\left[#2\right]}}
\newcommand{\Var}[1]{\mathify{{\rm VAR}\left[#1\right]}}
\newcommand{\Cov}{\mathify{{\rm Cov}}}
\newcommand{\Err}{\mathrm{Err}}
\newcommand{\Espace}{{\tilde{\mathcal{F}}[0,T]}}

% --- 5.8 Estimation / text-style operators
\newcommand{\tsnr}{{\text{snr}}}
\newcommand{\tvar}{{\text{Var}}}
\newcommand{\cmle}{\text{\rm cmle}}
\newcommand{\cmse}{\text{\rm cmse}}
\newcommand{\mmse}{\text{\rm mmse}}
\newcommand{\mmle}{\text{\rm mmle}}
\newcommand{\minimax}{\text{\rm minimax}}
\newcommand{\mse}{\text{\rm mse}}
\newcommand{\mle}{\text{\rm mle}}
\newcommand{\diag}{\mbox{diag}}
\newcommand{\JSD}{\mbox{JSD}}
\newcommand{\conv}{\text{\rm conv}}
\newcommand{\sgn}{{\mbox{sign}}}

\newcommand\argmin{\mathop{\mbox{{\rm argmin}}}\limits}
\newcommand\argmax{\mathop{\mbox{{\rm argmax}}}\limits}
\newcommand\maximize{\mathop{\mbox{{\rm maximize}}}\limits}
\newcommand\minimize{\mathop{\mbox{{\rm minimize}}}\limits}

% --- 5.9 Delimiters and relations
\newcommand{\lp}{\left(}
\newcommand{\rp}{\right)}
\newcommand{\lb}{\left[}
\newcommand{\rb}{\right]}
\newcommand{\lbp}{\left\{}
\newcommand{\rbp}{\right\}}
\newcommand{\lba}{\left\lvert}
\newcommand{\rba}{\right\rvert}
\newcommand{\lV}{\left\lVert}
\newcommand{\rV}{\right\rVert}
\newcommand{\mv}{\middle\vert}
\newcommand{\lan}{\langle}
\newcommand{\ran}{\rangle}
\newcommand{\ra}{\rightarrow}
\newcommand{\eqDef}{\coloneqq}   % needs mathtools
\newcommand{\suchthat}{\;\ifnum\currentgrouptype=16 \middle\fi|\;}
\newcommand{\dbar}{\;\ifnum\currentgrouptype=16 \middle \fi\|\;}
\newcommand{\equalto}[2]{\underset{\scriptstyle\overset{\mkern4mu\verteq}{#2}}{#1}}
\newcommand{\verteq}{\rotatebox{90}{$\,=$}}

% --- 5.10 Equation-array shortcuts
\newcommand{\bea}{\begin{eqnarray}}
\newcommand{\eea}{\end{eqnarray}}
\newcommand{\beas}{\begin{eqnarray*}}
\newcommand{\eeas}{\end{eqnarray*}}

% --- 5.11 Comment notes
\newcommand{\noa}[1]{{\color{red}\textbf{AN}: #1}}
\newcommand{\berivan}[1]{{\color{orange}\textbf{BI}: #1}}
\newcommand{\wn}[1]{{\color{green}\textbf{WN}: #1}}

% --- 5.12 Compatibility
% make hyperref and algorithmic work together better
\newcommand{\theHalgorithm}{\arabic{algorithm}}


%% =====================================================================
\begin{document}
%% =====================================================================

%% ---------------------------------------------------------------------
%%  Cover
%% ---------------------------------------------------------------------
\begin{center}
{\bf Introduction to Mathematics for Deep Learning: Midterm \begin{solution}Solutions\end{solution}} \\ \medskip
\begin{problem}Tuesday, October 20, 2026\end{problem}
\end{center}

\begin{problem}
\noindent
\begin{itemize}
\item You can use previous parts of a problem even if you did not solve them.
\item As throughout the course, entropy ($H$) and Mutual Information ($I$) are in bits.
\item $\log$ is taken in base 2; $\ln$ is the natural logarithm.
\end{itemize}
\LeftFoot{Midterm}
\end{problem}

\begin{solution}
\LeftFoot{Midterm Solution}
\end{solution}

\begin{problem}
    \vskip .3in
    {\bf ID : \underline{\hspace{6cm}}}\\
    \vskip .3in
    {\bf Name : \underline{\hspace{5.5cm}}}\\
    ~
    \vskip 2in
    \begin{table}[h]
        \Large
        \begin{tabular}{c|c}
           1     & \textcolor{white}{aaaaaaaa}/~40 \\
           2     & \textcolor{white}{aaaaaaaa}/~30 \\
           3     & \textcolor{white}{aaaaaaaa}/~30 \\
           \hline
           Total & \textcolor{white}{aaaaaaaa}/~100 \\
        \end{tabular}
    \end{table}
    \newpage
    {\bf extra sheet}
    \newpage
\end{problem}

\begin{enumerate}

%% =====================================================================
%%  Problem 1
%% =====================================================================
\item {\bf Mix of Questions} {\it (40 points, 10 points each)}

\begin{enumerate}[a)]

% (a)
\item Let $X_1,\dots,X_n$ be i.i.d.\ $\Bern(\frac12)$ random variables, $X^n=(X_1,\dots,X_n)$, and let $K=\sum_{i=1}^n X_i$ be the number of ones. Show that
\[
H(X^n\mid K)\ \ge\ n-\log(n+1).
\]

% (b)
\item Let $X$ and $Y$ be continuous random variables with densities $f_X$ and $f_Y$, respectively, and let $Z$ have density $f_Z=\frac12 f_X+\frac12 f_Y$. Prove or disprove: $h(Z)\le\max\{h(X),\,h(Y)\}$.

% (c)
\item Let $X$ be uniform on $K\ge 2$ classes $\{1,\dots,K\}$. A noisy label $Y$ is generated as follows: with probability $1-\varepsilon$, $Y=X$; with probability $\varepsilon$, $Y$ is drawn uniformly from the remaining $K-1$ classes (all classes other than $X$).
\begin{enumerate}[i)]
\item Compute $I(X;Y)$ for $\varepsilon=0$ and for $\varepsilon=1$, and show that both are nonzero.
\item Find $\varepsilon$ such that $I(X;Y)=0$.
\end{enumerate}

% (d)
\item A deck of $n$ cards is shuffled by the following move: take the top card and reinsert it at one of the $n$ possible positions (on top, below the 1st card, below the 2nd card, $\dots$, at the bottom), chosen uniformly at random and independently of all previous moves. Reinserting at the top means the card is put back where it was, so the deck stays the same. Let $X_0$ be the initial ordering of the deck and $X_i$ the ordering of the $n$ cards after the $i$-th move. Then $X_0,X_1,X_2,\dots$ is a Markov chain on the $n!$ orderings of the deck. Show that the uniform distribution over all $n!$ orderings is a stationary distribution of this Markov chain.
\end{enumerate}

%%% Solution of Problem 1
\begin{solution}
\vskip 0.2in
{\bf Solution:}

\begin{enumerate}[a)]

%% (a) sol
\item Since $K$ is a function of $X^n$, $H(X^n,K)=H(X^n)$. By the chain rule,
\[
n=H(X^n)=H(X^n,K)=H(K)+H(X^n\mid K).
\]
$K$ takes $n+1$ values $\{0,1,\dots,n\}$, so $H(K)\le\log(n+1)$. Hence
\[
H(X^n\mid K)=n-H(K)\ \ge\ n-\log(n+1).
\]
(Remark: the count $K$ carries at most $\log(n+1)$ of the $n$ bits; almost all of the information is in the order of the bits.)

%% (b) sol
\item The statement is {\bf false}. Counterexample: $X\sim{\rm Unif}[0,1]$, $Y\sim{\rm Unif}[10,11]$. Then
\[
h(X)=-\int_0^1 1\cdot\log 1\,dx=0,\qquad h(Y)=0.
\]
$Z$ has density $f_Z(z)=\frac12$ on $[0,1]\cup[10,11]$ (total length 2) and $0$ elsewhere, so
\[
h(Z)=-\int f_Z\log f_Z=-2\cdot\tfrac12\log\tfrac12=1\ >\ 0=\max\{h(X),h(Y)\}.
\]
(Remark: mixing adds the uncertainty about \emph{which} density $Z$ came from. Introducing a label $W\sim\Bern(\frac12)$ with $Z\mid\{W=1\}\sim f_X$ and $Z\mid\{W=0\}\sim f_Y$, one can show that, whenever $h(X)$ and $h(Y)$ are finite,
\[
\tfrac12h(X)+\tfrac12h(Y)\ \le\ h(Z)\ \le\ \tfrac12h(X)+\tfrac12h(Y)+1,
\]
with the upper bound attained exactly when $f_X(z)f_Y(z)=0$ for almost every $z$, as above.)

%% (c) sol
\item The channel is $P(Y=y\mid X=x)=1-\varepsilon$ if $y=x$ and $\frac{\varepsilon}{K-1}$ if $y\ne x$. By symmetry $Y$ is uniform, so $H(Y)=\log K$ and $H(X)=\log K$.
\begin{enumerate}[i)]
\item $\varepsilon=0$: $Y=X$, so $I(X;Y)=H(X)-H(X\mid Y)=\log K-0=\log K>0$.

$\varepsilon=1$: $Y$ is always wrong. Given $Y=y$, $X$ is uniform over the $K-1$ classes other than $y$, so $H(X\mid Y)=\log(K-1)$ and
\[
I(X;Y)=\log K-\log(K-1)=\log\frac{K}{K-1}>0.
\]
A label that is always wrong still tells us which class $X$ is \emph{not}. (For $K=2$ this is a full bit, $Y=3-X$; for $K\ge 3$ it is less than $\log K$.) Mutual information measures dependence, not correctness.

\item $I(X;Y)=0$ iff $X$ and $Y$ are independent, i.e., iff the conditional law $P(Y=\cdot\mid X=x)$ is the same for every $x$. That law puts $1-\varepsilon$ on $x$ and $\frac{\varepsilon}{K-1}$ on every other class, so it does not depend on $x$ iff it is uniform:
\[
1-\varepsilon=\frac{\varepsilon}{K-1}\iff\varepsilon=\frac{K-1}{K}.
\]
Then $Y$ is uniform regardless of $X$, so $I(X;Y)=0$. (Note: this is $\varepsilon=\frac12$ only for $K=2$; for $K=10$, labels that are wrong $60\%$ of the time are still informative.)
\end{enumerate}

%% (d) sol
\item Let $P(\tau\to\sigma)=P(X_{i+1}=\sigma\mid X_i=\tau)$ be the one-step transition probability between orderings and $U(\sigma)=\frac1{n!}$. We need $\sum_\tau U(\tau)P(\tau\to\sigma)=U(\sigma)$ for every ordering $\sigma$.
\begin{itemize}
\item From any ordering $\tau$, the $n$ insertion positions of the top card give $n$ \emph{different} orderings, each with probability $\frac1n$. So $P(\tau\to\sigma)\in\{0,\frac1n\}$.
\item Which $\tau$ can reach a fixed $\sigma$? Undo the move: pick any one of the $n$ cards of $\sigma$ and move it to the top. This gives exactly $n$ distinct predecessors $\tau$ (picking the top card of $\sigma$ gives $\tau=\sigma$, the ``reinsert on top'' move).
\end{itemize}
Therefore
\[
\sum_\tau U(\tau)P(\tau\to\sigma)=n\cdot\frac1{n!}\cdot\frac1n=\frac1{n!}=U(\sigma).
\]
(Equivalently, the transition matrix is doubly stochastic: every row and every column sums to 1, so the uniform distribution is stationary.)
\end{enumerate}
\end{solution}

\begin{problem}
\newpage {\bf problem 1 cont.}
\end{problem}
\newpage


%% =====================================================================
%%  Problem 2
%% =====================================================================
\item {\bf Label Smoothing vs.\ Confidence Penalty} {\it (30 points)}

A classifier for $K\ge 2$ classes outputs a probability distribution $q=(q_1,\dots,q_K)$ with all $q_k>0$, where $q_k$ is the probability it assigns to label $k$. Let $y$ be the true label, and let $u=(\frac1K,\dots,\frac1K)$ be the uniform distribution. Two common ways to discourage over-confident predictions are:
\begin{itemize}
\item {\bf Label smoothing.} For $0<\varepsilon<1$, replace the one-hot target by the \emph{smoothed one-hot} vector $o^{(\varepsilon)}=(o^{(\varepsilon)}_1,\dots,o^{(\varepsilon)}_K)$ with
\[
o^{(\varepsilon)}_k=
\begin{cases}
1-\varepsilon+\dfrac{\varepsilon}{K}, & k=y \ \text{(the true label)},\\[2mm]
\dfrac{\varepsilon}{K}, & k\ne y,
\end{cases}
\]
and minimize the cross-entropy $L_{\rm LS}(q)=-\sum_{k=1}^K o^{(\varepsilon)}_k\log q_k$. Note that $o^{(0)}$ is the one-hot vector of the true label $y$ (a $1$ in position $y$ and $0$ elsewhere), so $\varepsilon=0$ gives the usual cross-entropy loss $-\log q_y$.
\item {\bf Confidence penalty.} For $\beta>0$, minimize $L_{\rm CP}(q)=-\log q_y-\beta H(q)$, where $H(q)=-\sum_k q_k\log q_k$.
\end{itemize}

\begin{enumerate}[(a)]
\item Show the following: {\it (10 points)}
\begin{align*}
L_{\rm LS}(q)&=(1-\varepsilon)(-\log q_y)+\varepsilon\,D(u\|q)+\varepsilon\log K,\\
L_{\rm CP}(q)&=-\log q_y+\beta\,D(q\|u)-\beta\log K.
\end{align*}

\item Fix the true label $y$ and let the prediction become fully confident: $q_y\to 1$ (and of course $q_k\to 0$ for every $k\ne y$). Find the limits of $D(u\|q)$, $D(q\|u)$, $L_{\rm LS}(q)$ and $L_{\rm CP}(q)$. {\it (10 points)}

\item In this part only, $q$ may have zero entries. Let $p$ be any distribution on $\{1,\dots,K\}$.
\begin{enumerate}[i)]
\item Show that $D(p\|q)<\infty$ if and only if $q_k>0$ for every $k$ with $p_k>0$. {\it (4 points)}
\item Show that $D(q\|p)<\infty$ if and only if $q_k=0$ for every $k$ with $p_k=0$. {\it (4 points)}
\item Set $p=u$. For which $q$ is $D(u\|q)$ finite? For which $q$ is $D(q\|u)$ finite? {\it (2 points)}
\end{enumerate}
\end{enumerate}

%%% Solution of Problem 2
\begin{solution}
\vskip 0.2in
{\bf Solution:}

\begin{enumerate}[(a)]

%% (a) sol
\item Write $o^{(\varepsilon)}=(1-\varepsilon)o^{(0)}+\varepsilon u$, where $o^{(0)}$ is the one-hot vector of $y$. The cross-entropy is linear in the target, so
\[
L_{\rm LS}(q)=(1-\varepsilon)(-\log q_y)+\varepsilon\Big(-\sum_k\tfrac1K\log q_k\Big).
\]
For the second term,
\[
-\sum_k\tfrac1K\log q_k=\sum_k\tfrac1K\log\frac{1/K}{q_k}-\sum_k\tfrac1K\log\tfrac1K=D(u\|q)+\log K.
\]
This gives the first identity. For the second,
\[
D(q\|u)=\sum_k q_k\log\frac{q_k}{1/K}=\sum_k q_k\log q_k+\log K=\log K-H(q),
\]
so $-\beta H(q)=\beta D(q\|u)-\beta\log K$, which gives the second identity.

%% (b) sol
\item $D(u\|q)=\frac1K\sum_k\log\frac{1}{Kq_k}$. Each term with $k\ne y$ tends to $+\infty$ as $q_k\to 0$, so $D(u\|q)\to\infty$.

$D(q\|u)=\log K-H(q)\to\log K-0=\log K$, which is finite.

Hence $L_{\rm LS}(q)\to\infty$, since $(1-\varepsilon)(-\log q_y)\to 0$ but $\varepsilon D(u\|q)\to\infty$; and $L_{\rm CP}(q)\to 0+\beta\log K-\beta\log K=0$. Full confidence is infinitely expensive under label smoothing, but has finite loss under the confidence penalty.

%% (c) sol
\item Use the conventions $0\log\frac{0}{a}=0$ for $a\ge 0$ and $b\log\frac{b}{0}=+\infty$ for $b>0$.
\begin{enumerate}[i)]
\item $D(p\|q)=\sum_k p_k\log\frac{p_k}{q_k}$ is a finite sum. Terms with $p_k=0$ contribute $0$; a term with $p_k>0$ is finite if $q_k>0$ and $+\infty$ if $q_k=0$. Hence $D(p\|q)<\infty$ iff {\bf $q_k>0$ for every $k$ with $p_k>0$}.
\item $D(q\|p)=\sum_k q_k\log\frac{q_k}{p_k}$. Terms with $q_k=0$ contribute $0$; a term with $q_k>0$ is finite if $p_k>0$ and $+\infty$ if $p_k=0$. Hence $D(q\|p)<\infty$ iff {\bf $q_k=0$ for every $k$ with $p_k=0$}.
\item Since $u_k=\frac1K>0$ for every $k$:
\begin{itemize}
\item by (i), $D(u\|q)<\infty$ iff {\bf $q_k>0$ for every $k$} (and $D(u\|q)\to\infty$ as any $q_k\to 0$);
\item by (ii), the condition is empty, so $D(q\|u)<\infty$ for {\bf every} $q$; indeed $0\le D(q\|u)=\log K-H(q)\le\log K$.
\end{itemize}
\end{enumerate}
(Remark: this explains (b). $L_{\rm LS}$ contains $\varepsilon D(u\|q)$, which puts an unbounded cost on any class probability approaching $0$, so $L_{\rm LS}\to\infty$ as $q_y\to 1$. $L_{\rm CP}$ contains $\beta D(q\|u)$, which is bounded, so the fully confident limit has finite loss. Neither is a hard constraint: e.g.\ $L_{\rm LS}$ is minimized at $q=o^{(\varepsilon)}$, but any $q$ with all $q_k>0$ has finite loss. Minimizing $D(p\|q)$ over $q$ is called \emph{zero-avoiding} and minimizing $D(q\|p)$ \emph{zero-forcing}, after (i) and (ii).)
\end{enumerate}
\end{solution}

\begin{problem}
\newpage {\bf problem 2 cont.}
\end{problem}
\newpage


%% =====================================================================
%%  Problem 3
%% =====================================================================
\item {\bf Heavy Tails and the Median of Means} {\it (30 points)}

Let $X$ be a random variable with $X\ge 1$ and
\[
P(X>x)=x^{-3}\qquad\text{for all } x\ge 1.
\]
Let $X_1,\dots,X_n$ be i.i.d.\ copies of $X$, and $\bar X_n=\frac1n\sum_{i=1}^n X_i$. One can check that $\mathbb{E}[e^{tX}]=\infty$ for every $t>0$, so the Chernoff bound cannot be applied directly to $\bar X_n$.

{\it Hint.} For a nonnegative random variable $W$, $\mathbb{E}[W]=\int_0^\infty P(W>w)\,dw$. You may use this without proof.

\begin{enumerate}[(a)]
\item Show that $\mathbb{E}[X^2]=\int_0^\infty 2x\,P(X>x)\,dx$. Then compute $\mu=\mathbb{E}[X]$ and $\sigma^2={\rm Var}(X)$. {\it (5 points)}

\item For $a>0$, show that
\[
P(\bar X_n\ge\mu+a)\ \ge\ \big(n(\mu+a)\big)^{-3},
\]
and compare it with the upper bound given by Chebyshev's inequality. Conclude that $P(\bar X_n\ge\mu+a)$ decays polynomially, not exponentially, in $n$. [Hint: $\bar X_n\ge X_1/n$.] {\it (5 points)}

\item \emph{Median of means.} Split the $n$ samples into $k$ groups of $m$ samples each ($n=km$), compute the sample mean $\bar X^{(j)}$ of each group $j=1,\dots,k$, and let $\hat\mu$ be the median of the $k$ group means. Fix $a>0$ and call group $j$ \emph{bad} if $|\bar X^{(j)}-\mu|\ge a$. {\it (15 points)}
\begin{enumerate}[i)]
\item Show that if $m\ge 6/a^2$, each group is bad with probability at most $\frac18$. [Hint: Chebyshev's inequality for $\bar X^{(j)}$.] {\it (3 points)}
\item Show that if $|\hat\mu-\mu|\ge a$, then at least $k/2$ groups are bad. [Hint: if fewer than half of the groups are bad, where do more than half of the group means lie, and where is their median?] {\it (4 points)}
\item Let $Z$ be the number of bad groups. Since the groups use disjoint samples, you may assume the worst case $Z\sim{\rm Binomial}(k,\frac18)$. Apply the Chernoff bound to show that for every $t>0$,
\[
P(Z\ge k/2)\ \le\ \Big(e^{-t/2}\big(\tfrac78+\tfrac18 e^{t}\big)\Big)^k.
\]
[Hint: what is the moment generating function of a Bernoulli random variable?] {\it (4 points)}
\item Choose $e^t=7$ and show that the bound becomes $\big(\frac{\sqrt7}{4}\big)^k$. Combine (ii)--(iv) to conclude that with $m=\lceil 6/a^2\rceil$, $P(|\hat\mu-\mu|\ge a)$ decays exponentially in $n$. {\it (4 points)}
\end{enumerate}

\item Explain in two or three sentences why the median of means succeeds where the sample mean $\bar X_n$ fails. {\it (5 points)}
\end{enumerate}

%%% Solution of Problem 3
\begin{solution}
\vskip 0.2in
{\bf Solution:}

\begin{enumerate}[(a)]

%% (a) sol
\item Apply the hint to $W=X^2\ge 0$ and substitute $s=x^2$ ($ds=2x\,dx$):
\[
\mathbb{E}[X^2]=\int_0^\infty P(X^2>s)\,ds=\int_0^\infty P(X>\sqrt s)\,ds=\int_0^\infty 2x\,P(X>x)\,dx.
\]
Since $P(X>x)=1$ for $x<1$,
\[
\mu=1+\int_1^\infty x^{-3}\,dx=1+\tfrac12=\tfrac32,\qquad
\mathbb{E}[X^2]=1+\int_1^\infty 2x^{-2}\,dx=1+2=3,
\]
so $\sigma^2=3-\frac94=\frac34$. (Equivalently, use the density $f(x)=3x^{-4}$ for $x\ge 1$.)

%% (b) sol
\item All $X_i\ge 1>0$, so $\bar X_n\ge X_1/n$. If $X_1\ge n(\mu+a)$, then $\bar X_n\ge\mu+a$. Hence
\[
P(\bar X_n\ge\mu+a)\ \ge\ P\big(X_1\ge n(\mu+a)\big)=\big(n(\mu+a)\big)^{-3}.
\]
Chebyshev gives the upper bound
\[
P(\bar X_n\ge\mu+a)\le P(|\bar X_n-\mu|\ge a)\le\frac{\sigma^2}{na^2}=\frac{3}{4na^2}.
\]
So for fixed $a>0$ the probability tends to $0$ but is bounded above and below by polynomial rates ($n^{-1}$ and $n^{-3}$). In particular no bound $Ce^{-cn}$ with fixed $C,c>0$ can hold for all large $n$, because $(n(\mu+a))^{-3}\gg Ce^{-cn}$.

%% (c) sol
\item
\begin{enumerate}[i)]
\item $\bar X^{(j)}$ is the mean of $m$ i.i.d.\ copies of $X$, so ${\rm Var}(\bar X^{(j)})=\sigma^2/m$. By Chebyshev,
\[
P(|\bar X^{(j)}-\mu|\ge a)\le\frac{\sigma^2}{ma^2}=\frac{3}{4ma^2}\le\frac34\cdot\frac16=\frac18\qquad\text{when } m\ge 6/a^2.
\]
\item Suppose fewer than $k/2$ groups are bad. Then more than half of the group means lie in the interval $(\mu-a,\mu+a)$. If $\hat\mu\ge\mu+a$, then more than half of the group means would lie strictly below $\hat\mu$, contradicting the fact that a median has at least half of the values at or above it. The case $\hat\mu\le\mu-a$ is symmetric. Hence $|\hat\mu-\mu|<a$. Contrapositive: $|\hat\mu-\mu|\ge a$ implies at least $k/2$ bad groups, i.e., $\{|\hat\mu-\mu|\ge a\}\subseteq\{Z\ge k/2\}$.
\item Chernoff bound: for $t>0$, $e^{tz}$ is increasing, so by Markov's inequality
\[
P(Z\ge k/2)=P(e^{tZ}\ge e^{tk/2})\le e^{-tk/2}\,\mathbb{E}[e^{tZ}].
\]
$Z=\sum_{j=1}^k B_j$ with independent $B_j\sim\Bern(\frac18)$, and the MGF of a $\Bern(\frac18)$ variable is $\mathbb{E}[e^{tB_j}]=\frac78+\frac18e^t$. By independence, $\mathbb{E}[e^{tZ}]=\prod_j\mathbb{E}[e^{tB_j}]=\big(\frac78+\frac18e^t\big)^k$. Combining,
\[
P(Z\ge k/2)\le\Big(e^{-t/2}\big(\tfrac78+\tfrac18e^t\big)\Big)^k.
\]
\item With $e^t=7$: $e^{-t/2}\big(\frac78+\frac78\big)=\frac1{\sqrt7}\cdot\frac74=\frac{\sqrt7}{4}\approx 0.661<1$. By (ii) and (iii),
\[
P(|\hat\mu-\mu|\ge a)\le P(Z\ge k/2)\le\Big(\frac{\sqrt7}{4}\Big)^k.
\]
(If the bad-group probability is $q<\frac18$ rather than exactly $\frac18$, $Z$ is only smaller, so the bound still holds.) With $m=\lceil 6/a^2\rceil$ fixed, $k=n/m$ grows linearly in $n$, so the bound $\big(\frac{\sqrt7}{4}\big)^{n/m}$ decays exponentially in $n$. (The choice $e^t=7$ is the one that minimizes the bound in (iii).)
\end{enumerate}

%% (d) sol
\item A single sufficiently large sample can push $\bar X_n$ above $\mu+a$, and with a heavy tail such a sample appears with only polynomially small probability. The median of means confines any huge sample to a single group, where it spoils only one vote. The median fails only if half of the independent groups fail at once, which is a coin-flip large deviation with an exponentially small probability.
\end{enumerate}
\end{solution}

\begin{problem}
\newpage {\bf problem 3 cont.}
\end{problem}
\begin{problem}
\newpage {\bf extra sheet}
\newpage {\bf extra sheet}
\end{problem}

\end{enumerate}

\end{document}
"
%%  (run twice so "Page x of y" resolves).
%% =====================================================================

\documentclass[12pt]{article}

%% ---------------------------------------------------------------------
%%  0. Version switch: problem sheet vs. solutions
%% ---------------------------------------------------------------------
\newif\ifsolution
\solutionfalse            % \solutiontrue -> solutions version
\ifdefined\SOLUTION \solutiontrue \fi

%% ---------------------------------------------------------------------
%%  1. Page layout (A4)
%% ---------------------------------------------------------------------
\usepackage[a4paper,top=1in,bottom=1in,left=1in,right=1in,footskip=30pt]{geometry}

%% ---------------------------------------------------------------------
%%  2. Packages
%% ---------------------------------------------------------------------
% versions / lists / verbatim
\usepackage{versions}
\usepackage{enumerate}
%\usepackage{enumitem}
\usepackage{verbatim}

% math
\usepackage{amsmath}
\usepackage{amsthm}
\usepackage{amsopn}
\usepackage{amsfonts}

% graphics
\usepackage{graphicx}
\usepackage{tikz}
\usetikzlibrary{arrows,shapes}
\usetikzlibrary{fit}
\usepackage{pgfplots}
\usepackage{psfrag}
\usepackage{pstricks}

% header / footer
\usepackage{fancyhdr}
\usepackage{lastpage}

% references and notes (cleveref must come after hyperref)
\usepackage{hyperref}
\usepackage[capitalize,noabbrev]{cleveref}
% Todonotes is useful during development; to turn off comments use
%\usepackage[disable,textsize=tiny]{todonotes}
\usepackage[textsize=tiny]{todonotes}

\allowdisplaybreaks

%% ---------------------------------------------------------------------
%%  3. Header / footer
%% ---------------------------------------------------------------------
\pagestyle{fancy}
\renewcommand{\headrulewidth}{0pt}
\renewcommand{\footrulewidth}{0.3pt}
\newcommand{\hffont}{\fontfamily{phv}\fontsize{9}{11}\selectfont}
\def\hfsize{\footnotesize}
\def\pagenumber{Page \thepage\ of \pageref{LastPage}}
\fancyhf{}
\newcommand{\LeftFoot}[1]{\fancyfoot[LO]{\textsl{\hffont #1}}}
\fancyfoot[LE]{\textsl{\hffont \pagenumber}}
\fancyfoot[RE]{\textsl{\hffont Introduction to Mathematics for Deep Learning, 2026}}
\fancyfoot[RO]{\textsl{\hffont \pagenumber}}
%\thispagestyle{empty}

%% ---------------------------------------------------------------------
%%  4. Problem / solution environments
%% ---------------------------------------------------------------------
%  \begin{problem}...\end{problem}   : problem sheet only (cover, blank pages)
%  \begin{solution}...\end{solution} : solutions only
\ifsolution
  \excludeversion{problem}
  \includeversion{solution}
\else
  \includeversion{problem}
  \excludeversion{solution}
\fi

\renewenvironment{proof}{\noindent{\bf Proof}\hspace*{1em}}{\qed\bigskip\\}
\newcommand{\sol}{\noindent{\bf Solution: }}

%% ---------------------------------------------------------------------
%%  5. Macros
%% ---------------------------------------------------------------------

% --- 5.1 Calligraphic: \cA ... \cZ, \calD, \calG
\newcommand{\cA}{{\cal A}}
\newcommand{\cB}{{\cal B}}
\newcommand{\cC}{{\cal C}}
\newcommand{\cD}{{\cal D}}
\newcommand{\cE}{{\cal E}}
\newcommand{\cF}{{\cal F}}
\newcommand{\cG}{{\cal G}}
\newcommand{\cH}{{\cal H}}
\newcommand{\cI}{{\cal I}}
\newcommand{\cJ}{{\cal J}}
\newcommand{\cK}{{\cal K}}
\newcommand{\cL}{{\cal L}}
\newcommand{\cM}{{\cal M}}
\newcommand{\cN}{{\cal N}}
\newcommand{\cO}{{\cal O}}
\newcommand{\cP}{{\cal P}}
\newcommand{\cQ}{{\cal Q}}
\newcommand{\cR}{{\cal R}}
\newcommand{\cS}{{\cal S}}
\newcommand{\cT}{{\cal T}}
\newcommand{\cU}{{\cal U}}
\newcommand{\cV}{{\cal V}}
\newcommand{\cW}{{\cal W}}
\newcommand{\cX}{{\cal X}}
\newcommand{\cY}{{\cal Y}}
\newcommand{\cZ}{{\cal Z}}
\newcommand{\mcal}{\mathcal}

\DeclareMathAlphabet{\pazocal}{OMS}{zplm}{m}{n}
\newcommand{\calD}{{\pazocal {D}}}
\newcommand{\calG}{{\pazocal {G}}}

% --- 5.2 Bold lowercase: \ba ... \bz, \eb, \bmu
\newcommand{\ba}{{\bf a}}
\newcommand{\bb}{{\bf b}}
\newcommand{\bc}{{\bf c}}
\newcommand{\bd}{{\bf d}}
\newcommand{\be}{{\bf e}}
\newcommand{\eb}{{\bf e}}
%\newcommand{\bf}{{\bf f}}   % clashes with \bf
\newcommand{\bg}{{\bf g}}
\newcommand{\bh}{{\bf h}}
\newcommand{\bj}{{\bf j}}
\newcommand{\bk}{{\bf k}}
\newcommand{\bl}{{\bf l}}
\newcommand{\bm}{{\bf m}}
\newcommand{\bn}{{\bf n}}
\newcommand{\bo}{{\bf o}}
\newcommand{\bp}{{\bf p}}
\newcommand{\bq}{{\bf q}}
\newcommand{\br}{{\bf r}}
\newcommand{\bs}{{\bf s}}
\newcommand{\bt}{{\bf t}}
\newcommand{\bu}{{\bf u}}
\newcommand{\bv}{{\bf v}}
\newcommand{\bw}{{\bf w}}
\newcommand{\bx}{{\bf x}}
\newcommand{\by}{{\bf y}}
\newcommand{\bz}{{\bf z}}
\newcommand{\bmu}{{\pmb{\mu}}}

% --- 5.3 Bold uppercase: \bA ... \bZ
\newcommand{\bA}{{\bf A}}
\newcommand{\bB}{{\bf B}}
\newcommand{\bC}{{\bf C}}
\newcommand{\bD}{{\bf D}}
\newcommand{\bE}{{\bf E}}
\newcommand{\bF}{{\bf F}}
\newcommand{\bG}{{\bf G}}
\newcommand{\bH}{{\bf H}}
\newcommand{\bI}{{\bf I}}
\newcommand{\bJ}{{\bf J}}
\newcommand{\bK}{{\bf K}}
\newcommand{\bL}{{\bf L}}
\newcommand{\bM}{{\bf M}}
\newcommand{\bN}{{\bf N}}
\newcommand{\bO}{{\bf O}}
\newcommand{\bP}{{\bf P}}
\newcommand{\bQ}{{\bf Q}}
\newcommand{\bR}{{\bf R}}
\newcommand{\bS}{{\bf S}}
\newcommand{\bT}{{\bf T}}
\newcommand{\bU}{{\bf U}}
\newcommand{\bV}{{\bf V}}
\newcommand{\bW}{{\bf W}}
\newcommand{\bX}{{\bf X}}
\newcommand{\bY}{{\bf Y}}
\newcommand{\bZ}{{\bf Z}}

% --- 5.4 Bold vectors with decorations / superscripts
\newcommand{\bXt}{\tilde{\bf X}}
\newcommand{\bXh}{\hat{\bf X}}
\newcommand{\bNsup}[1]{{\bf N}^{(#1)}}
\newcommand{\bNtsup}[1]{\tilde{\bf N}^{(#1)}}
\newcommand{\bZsup}[1]{{\bf Z}^{(#1)}}
\newcommand{\bZtsup}[1]{\tilde{\bf Z}^{(#1)}}
\newcommand{\bXhsup}[1]{\hat{\bf X}^{(#1)}}
\newcommand{\bXtsup}[1]{\tilde{\bf X}^{(#1)}}
\newcommand{\bzerod}{{\bf 0}_d}
\newcommand{\boned}{{\bf 1}_d}
\newcommand{\bone}{{\bf 1}}

% indexed bold / plain symbols X^{(i)} etc.
\newcommand{\Xb}[1]{{\bf X^{(#1)}}}
\newcommand{\Xe}[1]{ X^{(#1)}}
\newcommand{\Bb}[1]{{\bf B^{(#1)}}}
\newcommand{\Be}[1]{ B^{(#1)}}
\newcommand{\Zb}[1]{{\bf Z^{(#1)}}}
\newcommand{\Ze}[1]{ Z^{(#1)}}
\newcommand{\Yb}[1]{{\bf Y^{(#1)}}}
\newcommand{\Ye}[1]{ Y^{(#1)}}
\newcommand{\Ub}[1]{{\bf U^{(#1)}}}
\newcommand{\Ue}[1]{ U^{(#1)}}
\newcommand{\Vb}[1]{{\bf V^{(#1)}}}
\newcommand{\Ve}[1]{ V^{(#1)}}
\newcommand{\Wb}[1]{{\bf W^{(#1)}}}
\newcommand{\We}[1]{ W^{(#1)}}
\newcommand{\Wbtilde}[1]{{\bf \widetilde{W}^{(#1)}}}
\newcommand{\Wetilde}[1]{ \widetilde{W}^{(#1)}}
\newcommand{\Xbh}[1]{{\bf \hat{X}^{(#1)}}}
\newcommand{\mb}[1]{{\bf m^{(#1)}}}
\newcommand{\me}[1]{ m^{(#1)}}
\newcommand{\mbhat}[1]{{\bf \hat{m}^{(#1)}}}
\newcommand{\mehat}[1]{{\hat{m}^{(#1)}}}
\newcommand{\Ab}{{\bf A}}
\newcommand{\Qb}{{\bf Q}}
\newcommand{\Xbhat}{{\bf \hat{X}}}

% --- 5.5 Tilde / hat / star
\newcommand{\tA}{\tilde{A}}
\newcommand{\tB}{\tilde{B}}
\newcommand{\tC}{\tilde{C}}
\newcommand{\tD}{\tilde{D}}
\newcommand{\tE}{\tilde{E}}
\newcommand{\tF}{\tilde{F}}
\newcommand{\tG}{\tilde{G}}
\newcommand{\tH}{\tilde{H}}
\newcommand{\tI}{\tilde{I}}
\newcommand{\tJ}{\tilde{J}}
\newcommand{\tK}{\tilde{K}}
\newcommand{\tL}{\tilde{L}}
\newcommand{\tM}{\tilde{M}}
\newcommand{\tN}{\tilde{N}}
\newcommand{\tO}{\tilde{O}}
\newcommand{\tP}{\tilde{P}}
\newcommand{\tQ}{\tilde{Q}}
\newcommand{\tR}{\tilde{R}}
\newcommand{\tS}{\tilde{S}}
\newcommand{\tT}{\tilde{T}}
\newcommand{\tU}{\tilde{U}}
\newcommand{\tV}{\tilde{V}}
\newcommand{\tX}{\tilde{X}}
\newcommand{\tY}{\tilde{Y}}
\newcommand{\tZ}{\tilde{Z}}
\newcommand{\tu}{\tilde{u}}
\newcommand{\tx}{\tilde{x}}
\newcommand{\ty}{\tilde{y}}
\newcommand{\tgamma}{\tilde{\gamma}}
%\newcommand{\ta}{\tilde{a}}

\newcommand{\uh}{\hat{u}}
\newcommand{\vh}{\hat{v}}
\newcommand{\xh}{\hat{x}}
\newcommand{\yh}{\hat{y}}
\newcommand{\zh}{\hat{z}}
\newcommand{\Ah}{\hat{A}}
\newcommand{\Uh}{\hat{U}}
\newcommand{\Xh}{\hat{X}}
\newcommand{\Yh}{\hat{Y}}
\newcommand{\Zh}{\hat{Z}}
\newcommand{\cXh}{\hat{\cX}}
\newcommand{\cYh}{\hat{\cY}}
\newcommand{\cZh}{\hat{\cZ}}
\newcommand{\bUh}{\hat{\bU}}

\newcommand{\bQstar}{\bQ^{\star}}
\newcommand{\Astar}{A^{\star}}
\newcommand{\Bstar}{B^{\star}}
\newcommand{\Qstar}{Q^{\star}}

% --- 5.6 Blackboard bold (number sets)
\newcommand{\field}[1]{\mathbb{#1}}
\newcommand{\fA}{\field{A}}
\newcommand{\fB}{\field{B}}
\newcommand{\fC}{\field{C}}
\newcommand{\fE}{\field{E}}
\newcommand{\fI}{\field{I}}
\newcommand{\fN}{\field{N}} % natural numbers
\newcommand{\fR}{\field{R}} % real numbers
\newcommand{\fS}{\field{S}}
\newcommand{\fT}{\field{T}}
\newcommand{\fZ}{\field{Z}} % integers

% --- 5.7 Probability: distributions, expectation, relations
\newcommand{\Bern}{{\mbox{Bern}}}
\newcommand{\Lap}{{\mbox{Lap}}}
\newcommand{\SpLap}{{\mbox{SpLap}}}
\newcommand{\punif}{P_{\text{unif}}}
\newcommand{\mathUO}{{\rm Unif}(\mathcal{O}(n))}
\newcommand{\iid}{\stackrel{iid}{\sim}}
\newcommand{\diid}{\overset{\text{i.i.d.}}{\sim}}
\newcommand{\deq}{\stackrel{(d)}{=}}

\newcommand{\mathify}[1]{\ifmmode{#1}\else\mbox{$#1$}\fi}
\newcommand{\Expec}[1]{\mathify{\mathbb{E}\left[#1\right]}}
\newcommand{\Esub}[2]{\mathify{\mathbb{E}_{#1}\left[#2\right]}}
% \newcommand{\Esubs}[2]{\mathify{\mathbb{E}_{#1}\left[#2\right]}}
\newcommand{\Var}[1]{\mathify{{\rm VAR}\left[#1\right]}}
\newcommand{\Cov}{\mathify{{\rm Cov}}}
\newcommand{\Err}{\mathrm{Err}}
\newcommand{\Espace}{{\tilde{\mathcal{F}}[0,T]}}

% --- 5.8 Estimation / text-style operators
\newcommand{\tsnr}{{\text{snr}}}
\newcommand{\tvar}{{\text{Var}}}
\newcommand{\cmle}{\text{\rm cmle}}
\newcommand{\cmse}{\text{\rm cmse}}
\newcommand{\mmse}{\text{\rm mmse}}
\newcommand{\mmle}{\text{\rm mmle}}
\newcommand{\minimax}{\text{\rm minimax}}
\newcommand{\mse}{\text{\rm mse}}
\newcommand{\mle}{\text{\rm mle}}
\newcommand{\diag}{\mbox{diag}}
\newcommand{\JSD}{\mbox{JSD}}
\newcommand{\conv}{\text{\rm conv}}
\newcommand{\sgn}{{\mbox{sign}}}

\newcommand\argmin{\mathop{\mbox{{\rm argmin}}}\limits}
\newcommand\argmax{\mathop{\mbox{{\rm argmax}}}\limits}
\newcommand\maximize{\mathop{\mbox{{\rm maximize}}}\limits}
\newcommand\minimize{\mathop{\mbox{{\rm minimize}}}\limits}

% --- 5.9 Delimiters and relations
\newcommand{\lp}{\left(}
\newcommand{\rp}{\right)}
\newcommand{\lb}{\left[}
\newcommand{\rb}{\right]}
\newcommand{\lbp}{\left\{}
\newcommand{\rbp}{\right\}}
\newcommand{\lba}{\left\lvert}
\newcommand{\rba}{\right\rvert}
\newcommand{\lV}{\left\lVert}
\newcommand{\rV}{\right\rVert}
\newcommand{\mv}{\middle\vert}
\newcommand{\lan}{\langle}
\newcommand{\ran}{\rangle}
\newcommand{\ra}{\rightarrow}
\newcommand{\eqDef}{\coloneqq}   % needs mathtools
\newcommand{\suchthat}{\;\ifnum\currentgrouptype=16 \middle\fi|\;}
\newcommand{\dbar}{\;\ifnum\currentgrouptype=16 \middle \fi\|\;}
\newcommand{\equalto}[2]{\underset{\scriptstyle\overset{\mkern4mu\verteq}{#2}}{#1}}
\newcommand{\verteq}{\rotatebox{90}{$\,=$}}

% --- 5.10 Equation-array shortcuts
\newcommand{\bea}{\begin{eqnarray}}
\newcommand{\eea}{\end{eqnarray}}
\newcommand{\beas}{\begin{eqnarray*}}
\newcommand{\eeas}{\end{eqnarray*}}

% --- 5.11 Comment notes
\newcommand{\noa}[1]{{\color{red}\textbf{AN}: #1}}
\newcommand{\berivan}[1]{{\color{orange}\textbf{BI}: #1}}
\newcommand{\wn}[1]{{\color{green}\textbf{WN}: #1}}

% --- 5.12 Compatibility
% make hyperref and algorithmic work together better
\newcommand{\theHalgorithm}{\arabic{algorithm}}


%% =====================================================================
\begin{document}
%% =====================================================================

%% ---------------------------------------------------------------------
%%  Cover
%% ---------------------------------------------------------------------
\begin{center}
{\bf Introduction to Mathematics for Deep Learning: Midterm \begin{solution}Solutions\end{solution}} \\ \medskip
\begin{problem}Tuesday, October 20, 2026\end{problem}
\end{center}

\begin{problem}
\noindent
\begin{itemize}
\item You can use previous parts of a problem even if you did not solve them.
\item As throughout the course, entropy ($H$) and Mutual Information ($I$) are in bits.
\item $\log$ is taken in base 2; $\ln$ is the natural logarithm.
\end{itemize}
\LeftFoot{Midterm}
\end{problem}

\begin{solution}
\LeftFoot{Midterm Solution}
\end{solution}

\begin{problem}
    \vskip .3in
    {\bf ID : \underline{\hspace{6cm}}}\\
    \vskip .3in
    {\bf Name : \underline{\hspace{5.5cm}}}\\
    ~
    \vskip 2in
    \begin{table}[h]
        \Large
        \begin{tabular}{c|c}
           1     & \textcolor{white}{aaaaaaaa}/~40 \\
           2     & \textcolor{white}{aaaaaaaa}/~30 \\
           3     & \textcolor{white}{aaaaaaaa}/~30 \\
           \hline
           Total & \textcolor{white}{aaaaaaaa}/~100 \\
        \end{tabular}
    \end{table}
    \newpage
    {\bf extra sheet}
    \newpage
\end{problem}

\begin{enumerate}

%% =====================================================================
%%  Problem 1
%% =====================================================================
\item {\bf Mix of Questions} {\it (40 points, 10 points each)}

\begin{enumerate}[a)]

% (a)
\item Let $X_1,\dots,X_n$ be i.i.d.\ $\Bern(\frac12)$ random variables, $X^n=(X_1,\dots,X_n)$, and let $K=\sum_{i=1}^n X_i$ be the number of ones. Show that
\[
H(X^n\mid K)\ \ge\ n-\log(n+1).
\]

% (b)
\item Let $X$ and $Y$ be continuous random variables with densities $f_X$ and $f_Y$, respectively, and let $Z$ have density $f_Z=\frac12 f_X+\frac12 f_Y$. Prove or disprove: $h(Z)\le\max\{h(X),\,h(Y)\}$.

% (c)
\item Let $X$ be uniform on $K\ge 2$ classes $\{1,\dots,K\}$. A noisy label $Y$ is generated as follows: with probability $1-\varepsilon$, $Y=X$; with probability $\varepsilon$, $Y$ is drawn uniformly from the remaining $K-1$ classes (all classes other than $X$).
\begin{enumerate}[i)]
\item Compute $I(X;Y)$ for $\varepsilon=0$ and for $\varepsilon=1$, and show that both are nonzero.
\item Find $\varepsilon$ such that $I(X;Y)=0$.
\end{enumerate}

% (d)
\item A deck of $n$ cards is shuffled by the following move: take the top card and reinsert it at one of the $n$ possible positions (on top, below the 1st card, below the 2nd card, $\dots$, at the bottom), chosen uniformly at random and independently of all previous moves. Reinserting at the top means the card is put back where it was, so the deck stays the same. Let $X_0$ be the initial ordering of the deck and $X_i$ the ordering of the $n$ cards after the $i$-th move. Then $X_0,X_1,X_2,\dots$ is a Markov chain on the $n!$ orderings of the deck. Show that the uniform distribution over all $n!$ orderings is a stationary distribution of this Markov chain.
\end{enumerate}

%%% Solution of Problem 1
\begin{solution}
\vskip 0.2in
{\bf Solution:}

\begin{enumerate}[a)]

%% (a) sol
\item Since $K$ is a function of $X^n$, $H(X^n,K)=H(X^n)$. By the chain rule,
\[
n=H(X^n)=H(X^n,K)=H(K)+H(X^n\mid K).
\]
$K$ takes $n+1$ values $\{0,1,\dots,n\}$, so $H(K)\le\log(n+1)$. Hence
\[
H(X^n\mid K)=n-H(K)\ \ge\ n-\log(n+1).
\]
(Remark: the count $K$ carries at most $\log(n+1)$ of the $n$ bits; almost all of the information is in the order of the bits.)

%% (b) sol
\item The statement is {\bf false}. Counterexample: $X\sim{\rm Unif}[0,1]$, $Y\sim{\rm Unif}[10,11]$. Then
\[
h(X)=-\int_0^1 1\cdot\log 1\,dx=0,\qquad h(Y)=0.
\]
$Z$ has density $f_Z(z)=\frac12$ on $[0,1]\cup[10,11]$ (total length 2) and $0$ elsewhere, so
\[
h(Z)=-\int f_Z\log f_Z=-2\cdot\tfrac12\log\tfrac12=1\ >\ 0=\max\{h(X),h(Y)\}.
\]
(Remark: mixing adds the uncertainty about \emph{which} density $Z$ came from. Introducing a label $W\sim\Bern(\frac12)$ with $Z\mid\{W=1\}\sim f_X$ and $Z\mid\{W=0\}\sim f_Y$, one can show that, whenever $h(X)$ and $h(Y)$ are finite,
\[
\tfrac12h(X)+\tfrac12h(Y)\ \le\ h(Z)\ \le\ \tfrac12h(X)+\tfrac12h(Y)+1,
\]
with the upper bound attained exactly when $f_X(z)f_Y(z)=0$ for almost every $z$, as above.)

%% (c) sol
\item The channel is $P(Y=y\mid X=x)=1-\varepsilon$ if $y=x$ and $\frac{\varepsilon}{K-1}$ if $y\ne x$. By symmetry $Y$ is uniform, so $H(Y)=\log K$ and $H(X)=\log K$.
\begin{enumerate}[i)]
\item $\varepsilon=0$: $Y=X$, so $I(X;Y)=H(X)-H(X\mid Y)=\log K-0=\log K>0$.

$\varepsilon=1$: $Y$ is always wrong. Given $Y=y$, $X$ is uniform over the $K-1$ classes other than $y$, so $H(X\mid Y)=\log(K-1)$ and
\[
I(X;Y)=\log K-\log(K-1)=\log\frac{K}{K-1}>0.
\]
A label that is always wrong still tells us which class $X$ is \emph{not}. (For $K=2$ this is a full bit, $Y=3-X$; for $K\ge 3$ it is less than $\log K$.) Mutual information measures dependence, not correctness.

\item $I(X;Y)=0$ iff $X$ and $Y$ are independent, i.e., iff the conditional law $P(Y=\cdot\mid X=x)$ is the same for every $x$. That law puts $1-\varepsilon$ on $x$ and $\frac{\varepsilon}{K-1}$ on every other class, so it does not depend on $x$ iff it is uniform:
\[
1-\varepsilon=\frac{\varepsilon}{K-1}\iff\varepsilon=\frac{K-1}{K}.
\]
Then $Y$ is uniform regardless of $X$, so $I(X;Y)=0$. (Note: this is $\varepsilon=\frac12$ only for $K=2$; for $K=10$, labels that are wrong $60\%$ of the time are still informative.)
\end{enumerate}

%% (d) sol
\item Let $P(\tau\to\sigma)=P(X_{i+1}=\sigma\mid X_i=\tau)$ be the one-step transition probability between orderings and $U(\sigma)=\frac1{n!}$. We need $\sum_\tau U(\tau)P(\tau\to\sigma)=U(\sigma)$ for every ordering $\sigma$.
\begin{itemize}
\item From any ordering $\tau$, the $n$ insertion positions of the top card give $n$ \emph{different} orderings, each with probability $\frac1n$. So $P(\tau\to\sigma)\in\{0,\frac1n\}$.
\item Which $\tau$ can reach a fixed $\sigma$? Undo the move: pick any one of the $n$ cards of $\sigma$ and move it to the top. This gives exactly $n$ distinct predecessors $\tau$ (picking the top card of $\sigma$ gives $\tau=\sigma$, the ``reinsert on top'' move).
\end{itemize}
Therefore
\[
\sum_\tau U(\tau)P(\tau\to\sigma)=n\cdot\frac1{n!}\cdot\frac1n=\frac1{n!}=U(\sigma).
\]
(Equivalently, the transition matrix is doubly stochastic: every row and every column sums to 1, so the uniform distribution is stationary.)
\end{enumerate}
\end{solution}

\begin{problem}
\newpage {\bf problem 1 cont.}
\end{problem}
\newpage


%% =====================================================================
%%  Problem 2
%% =====================================================================
\item {\bf Label Smoothing vs.\ Confidence Penalty} {\it (30 points)}

A classifier for $K\ge 2$ classes outputs a probability distribution $q=(q_1,\dots,q_K)$ with all $q_k>0$, where $q_k$ is the probability it assigns to label $k$. Let $y$ be the true label, and let $u=(\frac1K,\dots,\frac1K)$ be the uniform distribution. Two common ways to discourage over-confident predictions are:
\begin{itemize}
\item {\bf Label smoothing.} For $0<\varepsilon<1$, replace the one-hot target by the \emph{smoothed one-hot} vector $o^{(\varepsilon)}=(o^{(\varepsilon)}_1,\dots,o^{(\varepsilon)}_K)$ with
\[
o^{(\varepsilon)}_k=
\begin{cases}
1-\varepsilon+\dfrac{\varepsilon}{K}, & k=y \ \text{(the true label)},\\[2mm]
\dfrac{\varepsilon}{K}, & k\ne y,
\end{cases}
\]
and minimize the cross-entropy $L_{\rm LS}(q)=-\sum_{k=1}^K o^{(\varepsilon)}_k\log q_k$. Note that $o^{(0)}$ is the one-hot vector of the true label $y$ (a $1$ in position $y$ and $0$ elsewhere), so $\varepsilon=0$ gives the usual cross-entropy loss $-\log q_y$.
\item {\bf Confidence penalty.} For $\beta>0$, minimize $L_{\rm CP}(q)=-\log q_y-\beta H(q)$, where $H(q)=-\sum_k q_k\log q_k$.
\end{itemize}

\begin{enumerate}[(a)]
\item Show the following: {\it (10 points)}
\begin{align*}
L_{\rm LS}(q)&=(1-\varepsilon)(-\log q_y)+\varepsilon\,D(u\|q)+\varepsilon\log K,\\
L_{\rm CP}(q)&=-\log q_y+\beta\,D(q\|u)-\beta\log K.
\end{align*}

\item Fix the true label $y$ and let the prediction become fully confident: $q_y\to 1$ (and of course $q_k\to 0$ for every $k\ne y$). Find the limits of $D(u\|q)$, $D(q\|u)$, $L_{\rm LS}(q)$ and $L_{\rm CP}(q)$. {\it (10 points)}

\item In this part only, $q$ may have zero entries. Let $p$ be any distribution on $\{1,\dots,K\}$.
\begin{enumerate}[i)]
\item Show that $D(p\|q)<\infty$ if and only if $q_k>0$ for every $k$ with $p_k>0$. {\it (4 points)}
\item Show that $D(q\|p)<\infty$ if and only if $q_k=0$ for every $k$ with $p_k=0$. {\it (4 points)}
\item Set $p=u$. For which $q$ is $D(u\|q)$ finite? For which $q$ is $D(q\|u)$ finite? {\it (2 points)}
\end{enumerate}
\end{enumerate}

%%% Solution of Problem 2
\begin{solution}
\vskip 0.2in
{\bf Solution:}

\begin{enumerate}[(a)]

%% (a) sol
\item Write $o^{(\varepsilon)}=(1-\varepsilon)o^{(0)}+\varepsilon u$, where $o^{(0)}$ is the one-hot vector of $y$. The cross-entropy is linear in the target, so
\[
L_{\rm LS}(q)=(1-\varepsilon)(-\log q_y)+\varepsilon\Big(-\sum_k\tfrac1K\log q_k\Big).
\]
For the second term,
\[
-\sum_k\tfrac1K\log q_k=\sum_k\tfrac1K\log\frac{1/K}{q_k}-\sum_k\tfrac1K\log\tfrac1K=D(u\|q)+\log K.
\]
This gives the first identity. For the second,
\[
D(q\|u)=\sum_k q_k\log\frac{q_k}{1/K}=\sum_k q_k\log q_k+\log K=\log K-H(q),
\]
so $-\beta H(q)=\beta D(q\|u)-\beta\log K$, which gives the second identity.

%% (b) sol
\item $D(u\|q)=\frac1K\sum_k\log\frac{1}{Kq_k}$. Each term with $k\ne y$ tends to $+\infty$ as $q_k\to 0$, so $D(u\|q)\to\infty$.

$D(q\|u)=\log K-H(q)\to\log K-0=\log K$, which is finite.

Hence $L_{\rm LS}(q)\to\infty$, since $(1-\varepsilon)(-\log q_y)\to 0$ but $\varepsilon D(u\|q)\to\infty$; and $L_{\rm CP}(q)\to 0+\beta\log K-\beta\log K=0$. Full confidence is infinitely expensive under label smoothing, but has finite loss under the confidence penalty.

%% (c) sol
\item Use the conventions $0\log\frac{0}{a}=0$ for $a\ge 0$ and $b\log\frac{b}{0}=+\infty$ for $b>0$.
\begin{enumerate}[i)]
\item $D(p\|q)=\sum_k p_k\log\frac{p_k}{q_k}$ is a finite sum. Terms with $p_k=0$ contribute $0$; a term with $p_k>0$ is finite if $q_k>0$ and $+\infty$ if $q_k=0$. Hence $D(p\|q)<\infty$ iff {\bf $q_k>0$ for every $k$ with $p_k>0$}.
\item $D(q\|p)=\sum_k q_k\log\frac{q_k}{p_k}$. Terms with $q_k=0$ contribute $0$; a term with $q_k>0$ is finite if $p_k>0$ and $+\infty$ if $p_k=0$. Hence $D(q\|p)<\infty$ iff {\bf $q_k=0$ for every $k$ with $p_k=0$}.
\item Since $u_k=\frac1K>0$ for every $k$:
\begin{itemize}
\item by (i), $D(u\|q)<\infty$ iff {\bf $q_k>0$ for every $k$} (and $D(u\|q)\to\infty$ as any $q_k\to 0$);
\item by (ii), the condition is empty, so $D(q\|u)<\infty$ for {\bf every} $q$; indeed $0\le D(q\|u)=\log K-H(q)\le\log K$.
\end{itemize}
\end{enumerate}
(Remark: this explains (b). $L_{\rm LS}$ contains $\varepsilon D(u\|q)$, which puts an unbounded cost on any class probability approaching $0$, so $L_{\rm LS}\to\infty$ as $q_y\to 1$. $L_{\rm CP}$ contains $\beta D(q\|u)$, which is bounded, so the fully confident limit has finite loss. Neither is a hard constraint: e.g.\ $L_{\rm LS}$ is minimized at $q=o^{(\varepsilon)}$, but any $q$ with all $q_k>0$ has finite loss. Minimizing $D(p\|q)$ over $q$ is called \emph{zero-avoiding} and minimizing $D(q\|p)$ \emph{zero-forcing}, after (i) and (ii).)
\end{enumerate}
\end{solution}

\begin{problem}
\newpage {\bf problem 2 cont.}
\end{problem}
\newpage


%% =====================================================================
%%  Problem 3
%% =====================================================================
\item {\bf Heavy Tails and the Median of Means} {\it (30 points)}

Let $X$ be a random variable with $X\ge 1$ and
\[
P(X>x)=x^{-3}\qquad\text{for all } x\ge 1.
\]
Let $X_1,\dots,X_n$ be i.i.d.\ copies of $X$, and $\bar X_n=\frac1n\sum_{i=1}^n X_i$. One can check that $\mathbb{E}[e^{tX}]=\infty$ for every $t>0$, so the Chernoff bound cannot be applied directly to $\bar X_n$.

{\it Hint.} For a nonnegative random variable $W$, $\mathbb{E}[W]=\int_0^\infty P(W>w)\,dw$. You may use this without proof.

\begin{enumerate}[(a)]
\item Show that $\mathbb{E}[X^2]=\int_0^\infty 2x\,P(X>x)\,dx$. Then compute $\mu=\mathbb{E}[X]$ and $\sigma^2={\rm Var}(X)$. {\it (5 points)}

\item For $a>0$, show that
\[
P(\bar X_n\ge\mu+a)\ \ge\ \big(n(\mu+a)\big)^{-3},
\]
and compare it with the upper bound given by Chebyshev's inequality. Conclude that $P(\bar X_n\ge\mu+a)$ decays polynomially, not exponentially, in $n$. [Hint: $\bar X_n\ge X_1/n$.] {\it (5 points)}

\item \emph{Median of means.} Split the $n$ samples into $k$ groups of $m$ samples each ($n=km$), compute the sample mean $\bar X^{(j)}$ of each group $j=1,\dots,k$, and let $\hat\mu$ be the median of the $k$ group means. Fix $a>0$ and call group $j$ \emph{bad} if $|\bar X^{(j)}-\mu|\ge a$. {\it (15 points)}
\begin{enumerate}[i)]
\item Show that if $m\ge 6/a^2$, each group is bad with probability at most $\frac18$. [Hint: Chebyshev's inequality for $\bar X^{(j)}$.] {\it (3 points)}
\item Show that if $|\hat\mu-\mu|\ge a$, then at least $k/2$ groups are bad. [Hint: if fewer than half of the groups are bad, where do more than half of the group means lie, and where is their median?] {\it (4 points)}
\item Let $Z$ be the number of bad groups. Since the groups use disjoint samples, you may assume the worst case $Z\sim{\rm Binomial}(k,\frac18)$. Apply the Chernoff bound to show that for every $t>0$,
\[
P(Z\ge k/2)\ \le\ \Big(e^{-t/2}\big(\tfrac78+\tfrac18 e^{t}\big)\Big)^k.
\]
[Hint: what is the moment generating function of a Bernoulli random variable?] {\it (4 points)}
\item Choose $e^t=7$ and show that the bound becomes $\big(\frac{\sqrt7}{4}\big)^k$. Combine (ii)--(iv) to conclude that with $m=\lceil 6/a^2\rceil$, $P(|\hat\mu-\mu|\ge a)$ decays exponentially in $n$. {\it (4 points)}
\end{enumerate}

\item Explain in two or three sentences why the median of means succeeds where the sample mean $\bar X_n$ fails. {\it (5 points)}
\end{enumerate}

%%% Solution of Problem 3
\begin{solution}
\vskip 0.2in
{\bf Solution:}

\begin{enumerate}[(a)]

%% (a) sol
\item Apply the hint to $W=X^2\ge 0$ and substitute $s=x^2$ ($ds=2x\,dx$):
\[
\mathbb{E}[X^2]=\int_0^\infty P(X^2>s)\,ds=\int_0^\infty P(X>\sqrt s)\,ds=\int_0^\infty 2x\,P(X>x)\,dx.
\]
Since $P(X>x)=1$ for $x<1$,
\[
\mu=1+\int_1^\infty x^{-3}\,dx=1+\tfrac12=\tfrac32,\qquad
\mathbb{E}[X^2]=1+\int_1^\infty 2x^{-2}\,dx=1+2=3,
\]
so $\sigma^2=3-\frac94=\frac34$. (Equivalently, use the density $f(x)=3x^{-4}$ for $x\ge 1$.)

%% (b) sol
\item All $X_i\ge 1>0$, so $\bar X_n\ge X_1/n$. If $X_1\ge n(\mu+a)$, then $\bar X_n\ge\mu+a$. Hence
\[
P(\bar X_n\ge\mu+a)\ \ge\ P\big(X_1\ge n(\mu+a)\big)=\big(n(\mu+a)\big)^{-3}.
\]
Chebyshev gives the upper bound
\[
P(\bar X_n\ge\mu+a)\le P(|\bar X_n-\mu|\ge a)\le\frac{\sigma^2}{na^2}=\frac{3}{4na^2}.
\]
So for fixed $a>0$ the probability tends to $0$ but is bounded above and below by polynomial rates ($n^{-1}$ and $n^{-3}$). In particular no bound $Ce^{-cn}$ with fixed $C,c>0$ can hold for all large $n$, because $(n(\mu+a))^{-3}\gg Ce^{-cn}$.

%% (c) sol
\item
\begin{enumerate}[i)]
\item $\bar X^{(j)}$ is the mean of $m$ i.i.d.\ copies of $X$, so ${\rm Var}(\bar X^{(j)})=\sigma^2/m$. By Chebyshev,
\[
P(|\bar X^{(j)}-\mu|\ge a)\le\frac{\sigma^2}{ma^2}=\frac{3}{4ma^2}\le\frac34\cdot\frac16=\frac18\qquad\text{when } m\ge 6/a^2.
\]
\item Suppose fewer than $k/2$ groups are bad. Then more than half of the group means lie in the interval $(\mu-a,\mu+a)$. If $\hat\mu\ge\mu+a$, then more than half of the group means would lie strictly below $\hat\mu$, contradicting the fact that a median has at least half of the values at or above it. The case $\hat\mu\le\mu-a$ is symmetric. Hence $|\hat\mu-\mu|<a$. Contrapositive: $|\hat\mu-\mu|\ge a$ implies at least $k/2$ bad groups, i.e., $\{|\hat\mu-\mu|\ge a\}\subseteq\{Z\ge k/2\}$.
\item Chernoff bound: for $t>0$, $e^{tz}$ is increasing, so by Markov's inequality
\[
P(Z\ge k/2)=P(e^{tZ}\ge e^{tk/2})\le e^{-tk/2}\,\mathbb{E}[e^{tZ}].
\]
$Z=\sum_{j=1}^k B_j$ with independent $B_j\sim\Bern(\frac18)$, and the MGF of a $\Bern(\frac18)$ variable is $\mathbb{E}[e^{tB_j}]=\frac78+\frac18e^t$. By independence, $\mathbb{E}[e^{tZ}]=\prod_j\mathbb{E}[e^{tB_j}]=\big(\frac78+\frac18e^t\big)^k$. Combining,
\[
P(Z\ge k/2)\le\Big(e^{-t/2}\big(\tfrac78+\tfrac18e^t\big)\Big)^k.
\]
\item With $e^t=7$: $e^{-t/2}\big(\frac78+\frac78\big)=\frac1{\sqrt7}\cdot\frac74=\frac{\sqrt7}{4}\approx 0.661<1$. By (ii) and (iii),
\[
P(|\hat\mu-\mu|\ge a)\le P(Z\ge k/2)\le\Big(\frac{\sqrt7}{4}\Big)^k.
\]
(If the bad-group probability is $q<\frac18$ rather than exactly $\frac18$, $Z$ is only smaller, so the bound still holds.) With $m=\lceil 6/a^2\rceil$ fixed, $k=n/m$ grows linearly in $n$, so the bound $\big(\frac{\sqrt7}{4}\big)^{n/m}$ decays exponentially in $n$. (The choice $e^t=7$ is the one that minimizes the bound in (iii).)
\end{enumerate}

%% (d) sol
\item A single sufficiently large sample can push $\bar X_n$ above $\mu+a$, and with a heavy tail such a sample appears with only polynomially small probability. The median of means confines any huge sample to a single group, where it spoils only one vote. The median fails only if half of the independent groups fail at once, which is a coin-flip large deviation with an exponentially small probability.
\end{enumerate}
\end{solution}

\begin{problem}
\newpage {\bf problem 3 cont.}
\end{problem}
\begin{problem}
\newpage {\bf extra sheet}
\newpage {\bf extra sheet}
\end{problem}

\end{enumerate}

\end{document}
"
%%  (run twice so "Page x of y" resolves).
%% =====================================================================

\documentclass[12pt]{article}

%% ---------------------------------------------------------------------
%%  0. Version switch: problem sheet vs. solutions
%% ---------------------------------------------------------------------
\newif\ifsolution
\solutionfalse            % \solutiontrue -> solutions version
\ifdefined\SOLUTION \solutiontrue \fi

%% ---------------------------------------------------------------------
%%  1. Page layout (A4)
%% ---------------------------------------------------------------------
\usepackage[a4paper,top=1in,bottom=1in,left=1in,right=1in,footskip=30pt]{geometry}

%% ---------------------------------------------------------------------
%%  2. Packages
%% ---------------------------------------------------------------------
% versions / lists / verbatim
\usepackage{versions}
\usepackage{enumerate}
%\usepackage{enumitem}
\usepackage{verbatim}

% math
\usepackage{amsmath}
\usepackage{amsthm}
\usepackage{amsopn}
\usepackage{amsfonts}

% graphics
\usepackage{graphicx}
\usepackage{tikz}
\usetikzlibrary{arrows,shapes}
\usetikzlibrary{fit}
\usepackage{pgfplots}
\usepackage{psfrag}
\usepackage{pstricks}

% header / footer
\usepackage{fancyhdr}
\usepackage{lastpage}

% references and notes (cleveref must come after hyperref)
\usepackage{hyperref}
\usepackage[capitalize,noabbrev]{cleveref}
% Todonotes is useful during development; to turn off comments use
%\usepackage[disable,textsize=tiny]{todonotes}
\usepackage[textsize=tiny]{todonotes}

\allowdisplaybreaks

%% ---------------------------------------------------------------------
%%  3. Header / footer
%% ---------------------------------------------------------------------
\pagestyle{fancy}
\renewcommand{\headrulewidth}{0pt}
\renewcommand{\footrulewidth}{0.3pt}
\newcommand{\hffont}{\fontfamily{phv}\fontsize{9}{11}\selectfont}
\def\hfsize{\footnotesize}
\def\pagenumber{Page \thepage\ of \pageref{LastPage}}
\fancyhf{}
\newcommand{\LeftFoot}[1]{\fancyfoot[LO]{\textsl{\hffont #1}}}
\fancyfoot[LE]{\textsl{\hffont \pagenumber}}
\fancyfoot[RE]{\textsl{\hffont Introduction to Mathematics for Deep Learning, 2026}}
\fancyfoot[RO]{\textsl{\hffont \pagenumber}}
%\thispagestyle{empty}

%% ---------------------------------------------------------------------
%%  4. Problem / solution environments
%% ---------------------------------------------------------------------
%  \begin{problem}...\end{problem}   : problem sheet only (cover, blank pages)
%  \begin{solution}...\end{solution} : solutions only
\ifsolution
  \excludeversion{problem}
  \includeversion{solution}
\else
  \includeversion{problem}
  \excludeversion{solution}
\fi

\renewenvironment{proof}{\noindent{\bf Proof}\hspace*{1em}}{\qed\bigskip\\}
\newcommand{\sol}{\noindent{\bf Solution: }}

%% ---------------------------------------------------------------------
%%  5. Macros
%% ---------------------------------------------------------------------

% --- 5.1 Calligraphic: \cA ... \cZ, \calD, \calG
\newcommand{\cA}{{\cal A}}
\newcommand{\cB}{{\cal B}}
\newcommand{\cC}{{\cal C}}
\newcommand{\cD}{{\cal D}}
\newcommand{\cE}{{\cal E}}
\newcommand{\cF}{{\cal F}}
\newcommand{\cG}{{\cal G}}
\newcommand{\cH}{{\cal H}}
\newcommand{\cI}{{\cal I}}
\newcommand{\cJ}{{\cal J}}
\newcommand{\cK}{{\cal K}}
\newcommand{\cL}{{\cal L}}
\newcommand{\cM}{{\cal M}}
\newcommand{\cN}{{\cal N}}
\newcommand{\cO}{{\cal O}}
\newcommand{\cP}{{\cal P}}
\newcommand{\cQ}{{\cal Q}}
\newcommand{\cR}{{\cal R}}
\newcommand{\cS}{{\cal S}}
\newcommand{\cT}{{\cal T}}
\newcommand{\cU}{{\cal U}}
\newcommand{\cV}{{\cal V}}
\newcommand{\cW}{{\cal W}}
\newcommand{\cX}{{\cal X}}
\newcommand{\cY}{{\cal Y}}
\newcommand{\cZ}{{\cal Z}}
\newcommand{\mcal}{\mathcal}

\DeclareMathAlphabet{\pazocal}{OMS}{zplm}{m}{n}
\newcommand{\calD}{{\pazocal {D}}}
\newcommand{\calG}{{\pazocal {G}}}

% --- 5.2 Bold lowercase: \ba ... \bz, \eb, \bmu
\newcommand{\ba}{{\bf a}}
\newcommand{\bb}{{\bf b}}
\newcommand{\bc}{{\bf c}}
\newcommand{\bd}{{\bf d}}
\newcommand{\be}{{\bf e}}
\newcommand{\eb}{{\bf e}}
%\newcommand{\bf}{{\bf f}}   % clashes with \bf
\newcommand{\bg}{{\bf g}}
\newcommand{\bh}{{\bf h}}
\newcommand{\bj}{{\bf j}}
\newcommand{\bk}{{\bf k}}
\newcommand{\bl}{{\bf l}}
\newcommand{\bm}{{\bf m}}
\newcommand{\bn}{{\bf n}}
\newcommand{\bo}{{\bf o}}
\newcommand{\bp}{{\bf p}}
\newcommand{\bq}{{\bf q}}
\newcommand{\br}{{\bf r}}
\newcommand{\bs}{{\bf s}}
\newcommand{\bt}{{\bf t}}
\newcommand{\bu}{{\bf u}}
\newcommand{\bv}{{\bf v}}
\newcommand{\bw}{{\bf w}}
\newcommand{\bx}{{\bf x}}
\newcommand{\by}{{\bf y}}
\newcommand{\bz}{{\bf z}}
\newcommand{\bmu}{{\pmb{\mu}}}

% --- 5.3 Bold uppercase: \bA ... \bZ
\newcommand{\bA}{{\bf A}}
\newcommand{\bB}{{\bf B}}
\newcommand{\bC}{{\bf C}}
\newcommand{\bD}{{\bf D}}
\newcommand{\bE}{{\bf E}}
\newcommand{\bF}{{\bf F}}
\newcommand{\bG}{{\bf G}}
\newcommand{\bH}{{\bf H}}
\newcommand{\bI}{{\bf I}}
\newcommand{\bJ}{{\bf J}}
\newcommand{\bK}{{\bf K}}
\newcommand{\bL}{{\bf L}}
\newcommand{\bM}{{\bf M}}
\newcommand{\bN}{{\bf N}}
\newcommand{\bO}{{\bf O}}
\newcommand{\bP}{{\bf P}}
\newcommand{\bQ}{{\bf Q}}
\newcommand{\bR}{{\bf R}}
\newcommand{\bS}{{\bf S}}
\newcommand{\bT}{{\bf T}}
\newcommand{\bU}{{\bf U}}
\newcommand{\bV}{{\bf V}}
\newcommand{\bW}{{\bf W}}
\newcommand{\bX}{{\bf X}}
\newcommand{\bY}{{\bf Y}}
\newcommand{\bZ}{{\bf Z}}

% --- 5.4 Bold vectors with decorations / superscripts
\newcommand{\bXt}{\tilde{\bf X}}
\newcommand{\bXh}{\hat{\bf X}}
\newcommand{\bNsup}[1]{{\bf N}^{(#1)}}
\newcommand{\bNtsup}[1]{\tilde{\bf N}^{(#1)}}
\newcommand{\bZsup}[1]{{\bf Z}^{(#1)}}
\newcommand{\bZtsup}[1]{\tilde{\bf Z}^{(#1)}}
\newcommand{\bXhsup}[1]{\hat{\bf X}^{(#1)}}
\newcommand{\bXtsup}[1]{\tilde{\bf X}^{(#1)}}
\newcommand{\bzerod}{{\bf 0}_d}
\newcommand{\boned}{{\bf 1}_d}
\newcommand{\bone}{{\bf 1}}

% indexed bold / plain symbols X^{(i)} etc.
\newcommand{\Xb}[1]{{\bf X^{(#1)}}}
\newcommand{\Xe}[1]{ X^{(#1)}}
\newcommand{\Bb}[1]{{\bf B^{(#1)}}}
\newcommand{\Be}[1]{ B^{(#1)}}
\newcommand{\Zb}[1]{{\bf Z^{(#1)}}}
\newcommand{\Ze}[1]{ Z^{(#1)}}
\newcommand{\Yb}[1]{{\bf Y^{(#1)}}}
\newcommand{\Ye}[1]{ Y^{(#1)}}
\newcommand{\Ub}[1]{{\bf U^{(#1)}}}
\newcommand{\Ue}[1]{ U^{(#1)}}
\newcommand{\Vb}[1]{{\bf V^{(#1)}}}
\newcommand{\Ve}[1]{ V^{(#1)}}
\newcommand{\Wb}[1]{{\bf W^{(#1)}}}
\newcommand{\We}[1]{ W^{(#1)}}
\newcommand{\Wbtilde}[1]{{\bf \widetilde{W}^{(#1)}}}
\newcommand{\Wetilde}[1]{ \widetilde{W}^{(#1)}}
\newcommand{\Xbh}[1]{{\bf \hat{X}^{(#1)}}}
\newcommand{\mb}[1]{{\bf m^{(#1)}}}
\newcommand{\me}[1]{ m^{(#1)}}
\newcommand{\mbhat}[1]{{\bf \hat{m}^{(#1)}}}
\newcommand{\mehat}[1]{{\hat{m}^{(#1)}}}
\newcommand{\Ab}{{\bf A}}
\newcommand{\Qb}{{\bf Q}}
\newcommand{\Xbhat}{{\bf \hat{X}}}

% --- 5.5 Tilde / hat / star
\newcommand{\tA}{\tilde{A}}
\newcommand{\tB}{\tilde{B}}
\newcommand{\tC}{\tilde{C}}
\newcommand{\tD}{\tilde{D}}
\newcommand{\tE}{\tilde{E}}
\newcommand{\tF}{\tilde{F}}
\newcommand{\tG}{\tilde{G}}
\newcommand{\tH}{\tilde{H}}
\newcommand{\tI}{\tilde{I}}
\newcommand{\tJ}{\tilde{J}}
\newcommand{\tK}{\tilde{K}}
\newcommand{\tL}{\tilde{L}}
\newcommand{\tM}{\tilde{M}}
\newcommand{\tN}{\tilde{N}}
\newcommand{\tO}{\tilde{O}}
\newcommand{\tP}{\tilde{P}}
\newcommand{\tQ}{\tilde{Q}}
\newcommand{\tR}{\tilde{R}}
\newcommand{\tS}{\tilde{S}}
\newcommand{\tT}{\tilde{T}}
\newcommand{\tU}{\tilde{U}}
\newcommand{\tV}{\tilde{V}}
\newcommand{\tX}{\tilde{X}}
\newcommand{\tY}{\tilde{Y}}
\newcommand{\tZ}{\tilde{Z}}
\newcommand{\tu}{\tilde{u}}
\newcommand{\tx}{\tilde{x}}
\newcommand{\ty}{\tilde{y}}
\newcommand{\tgamma}{\tilde{\gamma}}
%\newcommand{\ta}{\tilde{a}}

\newcommand{\uh}{\hat{u}}
\newcommand{\vh}{\hat{v}}
\newcommand{\xh}{\hat{x}}
\newcommand{\yh}{\hat{y}}
\newcommand{\zh}{\hat{z}}
\newcommand{\Ah}{\hat{A}}
\newcommand{\Uh}{\hat{U}}
\newcommand{\Xh}{\hat{X}}
\newcommand{\Yh}{\hat{Y}}
\newcommand{\Zh}{\hat{Z}}
\newcommand{\cXh}{\hat{\cX}}
\newcommand{\cYh}{\hat{\cY}}
\newcommand{\cZh}{\hat{\cZ}}
\newcommand{\bUh}{\hat{\bU}}

\newcommand{\bQstar}{\bQ^{\star}}
\newcommand{\Astar}{A^{\star}}
\newcommand{\Bstar}{B^{\star}}
\newcommand{\Qstar}{Q^{\star}}

% --- 5.6 Blackboard bold (number sets)
\newcommand{\field}[1]{\mathbb{#1}}
\newcommand{\fA}{\field{A}}
\newcommand{\fB}{\field{B}}
\newcommand{\fC}{\field{C}}
\newcommand{\fE}{\field{E}}
\newcommand{\fI}{\field{I}}
\newcommand{\fN}{\field{N}} % natural numbers
\newcommand{\fR}{\field{R}} % real numbers
\newcommand{\fS}{\field{S}}
\newcommand{\fT}{\field{T}}
\newcommand{\fZ}{\field{Z}} % integers

% --- 5.7 Probability: distributions, expectation, relations
\newcommand{\Bern}{{\mbox{Bern}}}
\newcommand{\Lap}{{\mbox{Lap}}}
\newcommand{\SpLap}{{\mbox{SpLap}}}
\newcommand{\punif}{P_{\text{unif}}}
\newcommand{\mathUO}{{\rm Unif}(\mathcal{O}(n))}
\newcommand{\iid}{\stackrel{iid}{\sim}}
\newcommand{\diid}{\overset{\text{i.i.d.}}{\sim}}
\newcommand{\deq}{\stackrel{(d)}{=}}

\newcommand{\mathify}[1]{\ifmmode{#1}\else\mbox{$#1$}\fi}
\newcommand{\Expec}[1]{\mathify{\mathbb{E}\left[#1\right]}}
\newcommand{\Esub}[2]{\mathify{\mathbb{E}_{#1}\left[#2\right]}}
% \newcommand{\Esubs}[2]{\mathify{\mathbb{E}_{#1}\left[#2\right]}}
\newcommand{\Var}[1]{\mathify{{\rm VAR}\left[#1\right]}}
\newcommand{\Cov}{\mathify{{\rm Cov}}}
\newcommand{\Err}{\mathrm{Err}}
\newcommand{\Espace}{{\tilde{\mathcal{F}}[0,T]}}

% --- 5.8 Estimation / text-style operators
\newcommand{\tsnr}{{\text{snr}}}
\newcommand{\tvar}{{\text{Var}}}
\newcommand{\cmle}{\text{\rm cmle}}
\newcommand{\cmse}{\text{\rm cmse}}
\newcommand{\mmse}{\text{\rm mmse}}
\newcommand{\mmle}{\text{\rm mmle}}
\newcommand{\minimax}{\text{\rm minimax}}
\newcommand{\mse}{\text{\rm mse}}
\newcommand{\mle}{\text{\rm mle}}
\newcommand{\diag}{\mbox{diag}}
\newcommand{\JSD}{\mbox{JSD}}
\newcommand{\conv}{\text{\rm conv}}
\newcommand{\sgn}{{\mbox{sign}}}

\newcommand\argmin{\mathop{\mbox{{\rm argmin}}}\limits}
\newcommand\argmax{\mathop{\mbox{{\rm argmax}}}\limits}
\newcommand\maximize{\mathop{\mbox{{\rm maximize}}}\limits}
\newcommand\minimize{\mathop{\mbox{{\rm minimize}}}\limits}

% --- 5.9 Delimiters and relations
\newcommand{\lp}{\left(}
\newcommand{\rp}{\right)}
\newcommand{\lb}{\left[}
\newcommand{\rb}{\right]}
\newcommand{\lbp}{\left\{}
\newcommand{\rbp}{\right\}}
\newcommand{\lba}{\left\lvert}
\newcommand{\rba}{\right\rvert}
\newcommand{\lV}{\left\lVert}
\newcommand{\rV}{\right\rVert}
\newcommand{\mv}{\middle\vert}
\newcommand{\lan}{\langle}
\newcommand{\ran}{\rangle}
\newcommand{\ra}{\rightarrow}
\newcommand{\eqDef}{\coloneqq}   % needs mathtools
\newcommand{\suchthat}{\;\ifnum\currentgrouptype=16 \middle\fi|\;}
\newcommand{\dbar}{\;\ifnum\currentgrouptype=16 \middle \fi\|\;}
\newcommand{\equalto}[2]{\underset{\scriptstyle\overset{\mkern4mu\verteq}{#2}}{#1}}
\newcommand{\verteq}{\rotatebox{90}{$\,=$}}

% --- 5.10 Equation-array shortcuts
\newcommand{\bea}{\begin{eqnarray}}
\newcommand{\eea}{\end{eqnarray}}
\newcommand{\beas}{\begin{eqnarray*}}
\newcommand{\eeas}{\end{eqnarray*}}

% --- 5.11 Comment notes
\newcommand{\noa}[1]{{\color{red}\textbf{AN}: #1}}
\newcommand{\berivan}[1]{{\color{orange}\textbf{BI}: #1}}
\newcommand{\wn}[1]{{\color{green}\textbf{WN}: #1}}

% --- 5.12 Compatibility
% make hyperref and algorithmic work together better
\newcommand{\theHalgorithm}{\arabic{algorithm}}


%% =====================================================================
\begin{document}
%% =====================================================================

%% ---------------------------------------------------------------------
%%  Cover
%% ---------------------------------------------------------------------
\begin{center}
{\bf Introduction to Mathematics for Deep Learning: Midterm \begin{solution}Solutions\end{solution}} \\ \medskip
\begin{problem}Tuesday, October 20, 2026\end{problem}
\end{center}

\begin{problem}
\noindent
\begin{itemize}
\item You can use previous parts of a problem even if you did not solve them.
\item As throughout the course, entropy ($H$) and Mutual Information ($I$) are in bits.
\item $\log$ is taken in base 2; $\ln$ is the natural logarithm.
\end{itemize}
\LeftFoot{Midterm}
\end{problem}

\begin{solution}
\LeftFoot{Midterm Solution}
\end{solution}

\begin{problem}
    \vskip .3in
    {\bf ID : \underline{\hspace{6cm}}}\\
    \vskip .3in
    {\bf Name : \underline{\hspace{5.5cm}}}\\
    ~
    \vskip 2in
    \begin{table}[h]
        \Large
        \begin{tabular}{c|c}
           1     & \textcolor{white}{aaaaaaaa}/~40 \\
           2     & \textcolor{white}{aaaaaaaa}/~30 \\
           3     & \textcolor{white}{aaaaaaaa}/~30 \\
           \hline
           Total & \textcolor{white}{aaaaaaaa}/~100 \\
        \end{tabular}
    \end{table}
    \newpage
    {\bf extra sheet}
    \newpage
\end{problem}

\begin{enumerate}

%% =====================================================================
%%  Problem 1
%% =====================================================================
\item {\bf Mix of Questions} {\it (40 points, 10 points each)}

\begin{enumerate}[a)]

% (a)
\item Let $X_1,\dots,X_n$ be i.i.d.\ $\Bern(\frac12)$ random variables, $X^n=(X_1,\dots,X_n)$, and let $K=\sum_{i=1}^n X_i$ be the number of ones. Show that
\[
H(X^n\mid K)\ \ge\ n-\log(n+1).
\]

% (b)
\item Let $X$ and $Y$ be continuous random variables with densities $f_X$ and $f_Y$, respectively, and let $Z$ have density $f_Z=\frac12 f_X+\frac12 f_Y$. Prove or disprove: $h(Z)\le\max\{h(X),\,h(Y)\}$.

% (c)
\item Let $X$ be uniform on $K\ge 2$ classes $\{1,\dots,K\}$. A noisy label $Y$ is generated as follows: with probability $1-\varepsilon$, $Y=X$; with probability $\varepsilon$, $Y$ is drawn uniformly from the remaining $K-1$ classes (all classes other than $X$).
\begin{enumerate}[i)]
\item Compute $I(X;Y)$ for $\varepsilon=0$ and for $\varepsilon=1$, and show that both are nonzero.
\item Find $\varepsilon$ such that $I(X;Y)=0$.
\end{enumerate}

% (d)
\item A deck of $n$ cards is shuffled by the following move: take the top card and reinsert it at one of the $n$ possible positions (on top, below the 1st card, below the 2nd card, $\dots$, at the bottom), chosen uniformly at random and independently of all previous moves. Reinserting at the top means the card is put back where it was, so the deck stays the same. Let $X_0$ be the initial ordering of the deck and $X_i$ the ordering of the $n$ cards after the $i$-th move. Then $X_0,X_1,X_2,\dots$ is a Markov chain on the $n!$ orderings of the deck. Show that the uniform distribution over all $n!$ orderings is a stationary distribution of this Markov chain.
\end{enumerate}

%%% Solution of Problem 1
\begin{solution}
\vskip 0.2in
{\bf Solution:}

\begin{enumerate}[a)]

%% (a) sol
\item Since $K$ is a function of $X^n$, $H(X^n,K)=H(X^n)$. By the chain rule,
\[
n=H(X^n)=H(X^n,K)=H(K)+H(X^n\mid K).
\]
$K$ takes $n+1$ values $\{0,1,\dots,n\}$, so $H(K)\le\log(n+1)$. Hence
\[
H(X^n\mid K)=n-H(K)\ \ge\ n-\log(n+1).
\]
(Remark: the count $K$ carries at most $\log(n+1)$ of the $n$ bits; almost all of the information is in the order of the bits.)

%% (b) sol
\item The statement is {\bf false}. Counterexample: $X\sim{\rm Unif}[0,1]$, $Y\sim{\rm Unif}[10,11]$. Then
\[
h(X)=-\int_0^1 1\cdot\log 1\,dx=0,\qquad h(Y)=0.
\]
$Z$ has density $f_Z(z)=\frac12$ on $[0,1]\cup[10,11]$ (total length 2) and $0$ elsewhere, so
\[
h(Z)=-\int f_Z\log f_Z=-2\cdot\tfrac12\log\tfrac12=1\ >\ 0=\max\{h(X),h(Y)\}.
\]
(Remark: mixing adds the uncertainty about \emph{which} density $Z$ came from. Introducing a label $W\sim\Bern(\frac12)$ with $Z\mid\{W=1\}\sim f_X$ and $Z\mid\{W=0\}\sim f_Y$, one can show that, whenever $h(X)$ and $h(Y)$ are finite,
\[
\tfrac12h(X)+\tfrac12h(Y)\ \le\ h(Z)\ \le\ \tfrac12h(X)+\tfrac12h(Y)+1,
\]
with the upper bound attained exactly when $f_X(z)f_Y(z)=0$ for almost every $z$, as above.)

%% (c) sol
\item The channel is $P(Y=y\mid X=x)=1-\varepsilon$ if $y=x$ and $\frac{\varepsilon}{K-1}$ if $y\ne x$. By symmetry $Y$ is uniform, so $H(Y)=\log K$ and $H(X)=\log K$.
\begin{enumerate}[i)]
\item $\varepsilon=0$: $Y=X$, so $I(X;Y)=H(X)-H(X\mid Y)=\log K-0=\log K>0$.

$\varepsilon=1$: $Y$ is always wrong. Given $Y=y$, $X$ is uniform over the $K-1$ classes other than $y$, so $H(X\mid Y)=\log(K-1)$ and
\[
I(X;Y)=\log K-\log(K-1)=\log\frac{K}{K-1}>0.
\]
A label that is always wrong still tells us which class $X$ is \emph{not}. (For $K=2$ this is a full bit, $Y=3-X$; for $K\ge 3$ it is less than $\log K$.) Mutual information measures dependence, not correctness.

\item $I(X;Y)=0$ iff $X$ and $Y$ are independent, i.e., iff the conditional law $P(Y=\cdot\mid X=x)$ is the same for every $x$. That law puts $1-\varepsilon$ on $x$ and $\frac{\varepsilon}{K-1}$ on every other class, so it does not depend on $x$ iff it is uniform:
\[
1-\varepsilon=\frac{\varepsilon}{K-1}\iff\varepsilon=\frac{K-1}{K}.
\]
Then $Y$ is uniform regardless of $X$, so $I(X;Y)=0$. (Note: this is $\varepsilon=\frac12$ only for $K=2$; for $K=10$, labels that are wrong $60\%$ of the time are still informative.)
\end{enumerate}

%% (d) sol
\item Let $P(\tau\to\sigma)=P(X_{i+1}=\sigma\mid X_i=\tau)$ be the one-step transition probability between orderings and $U(\sigma)=\frac1{n!}$. We need $\sum_\tau U(\tau)P(\tau\to\sigma)=U(\sigma)$ for every ordering $\sigma$.
\begin{itemize}
\item From any ordering $\tau$, the $n$ insertion positions of the top card give $n$ \emph{different} orderings, each with probability $\frac1n$. So $P(\tau\to\sigma)\in\{0,\frac1n\}$.
\item Which $\tau$ can reach a fixed $\sigma$? Undo the move: pick any one of the $n$ cards of $\sigma$ and move it to the top. This gives exactly $n$ distinct predecessors $\tau$ (picking the top card of $\sigma$ gives $\tau=\sigma$, the ``reinsert on top'' move).
\end{itemize}
Therefore
\[
\sum_\tau U(\tau)P(\tau\to\sigma)=n\cdot\frac1{n!}\cdot\frac1n=\frac1{n!}=U(\sigma).
\]
(Equivalently, the transition matrix is doubly stochastic: every row and every column sums to 1, so the uniform distribution is stationary.)
\end{enumerate}
\end{solution}

\begin{problem}
\newpage {\bf problem 1 cont.}
\end{problem}
\newpage


%% =====================================================================
%%  Problem 2
%% =====================================================================
\item {\bf Label Smoothing vs.\ Confidence Penalty} {\it (30 points)}

A classifier for $K\ge 2$ classes outputs a probability distribution $q=(q_1,\dots,q_K)$ with all $q_k>0$, where $q_k$ is the probability it assigns to label $k$. Let $y$ be the true label, and let $u=(\frac1K,\dots,\frac1K)$ be the uniform distribution. Two common ways to discourage over-confident predictions are:
\begin{itemize}
\item {\bf Label smoothing.} For $0<\varepsilon<1$, replace the one-hot target by the \emph{smoothed one-hot} vector $o^{(\varepsilon)}=(o^{(\varepsilon)}_1,\dots,o^{(\varepsilon)}_K)$ with
\[
o^{(\varepsilon)}_k=
\begin{cases}
1-\varepsilon+\dfrac{\varepsilon}{K}, & k=y \ \text{(the true label)},\\[2mm]
\dfrac{\varepsilon}{K}, & k\ne y,
\end{cases}
\]
and minimize the cross-entropy $L_{\rm LS}(q)=-\sum_{k=1}^K o^{(\varepsilon)}_k\log q_k$. Note that $o^{(0)}$ is the one-hot vector of the true label $y$ (a $1$ in position $y$ and $0$ elsewhere), so $\varepsilon=0$ gives the usual cross-entropy loss $-\log q_y$.
\item {\bf Confidence penalty.} For $\beta>0$, minimize $L_{\rm CP}(q)=-\log q_y-\beta H(q)$, where $H(q)=-\sum_k q_k\log q_k$.
\end{itemize}

\begin{enumerate}[(a)]
\item Show the following: {\it (10 points)}
\begin{align*}
L_{\rm LS}(q)&=(1-\varepsilon)(-\log q_y)+\varepsilon\,D(u\|q)+\varepsilon\log K,\\
L_{\rm CP}(q)&=-\log q_y+\beta\,D(q\|u)-\beta\log K.
\end{align*}

\item Fix the true label $y$ and let the prediction become fully confident: $q_y\to 1$ (and of course $q_k\to 0$ for every $k\ne y$). Find the limits of $D(u\|q)$, $D(q\|u)$, $L_{\rm LS}(q)$ and $L_{\rm CP}(q)$. {\it (10 points)}

\item In this part only, $q$ may have zero entries. Let $p$ be any distribution on $\{1,\dots,K\}$.
\begin{enumerate}[i)]
\item Show that $D(p\|q)<\infty$ if and only if $q_k>0$ for every $k$ with $p_k>0$. {\it (4 points)}
\item Show that $D(q\|p)<\infty$ if and only if $q_k=0$ for every $k$ with $p_k=0$. {\it (4 points)}
\item Set $p=u$. For which $q$ is $D(u\|q)$ finite? For which $q$ is $D(q\|u)$ finite? {\it (2 points)}
\end{enumerate}
\end{enumerate}

%%% Solution of Problem 2
\begin{solution}
\vskip 0.2in
{\bf Solution:}

\begin{enumerate}[(a)]

%% (a) sol
\item Write $o^{(\varepsilon)}=(1-\varepsilon)o^{(0)}+\varepsilon u$, where $o^{(0)}$ is the one-hot vector of $y$. The cross-entropy is linear in the target, so
\[
L_{\rm LS}(q)=(1-\varepsilon)(-\log q_y)+\varepsilon\Big(-\sum_k\tfrac1K\log q_k\Big).
\]
For the second term,
\[
-\sum_k\tfrac1K\log q_k=\sum_k\tfrac1K\log\frac{1/K}{q_k}-\sum_k\tfrac1K\log\tfrac1K=D(u\|q)+\log K.
\]
This gives the first identity. For the second,
\[
D(q\|u)=\sum_k q_k\log\frac{q_k}{1/K}=\sum_k q_k\log q_k+\log K=\log K-H(q),
\]
so $-\beta H(q)=\beta D(q\|u)-\beta\log K$, which gives the second identity.

%% (b) sol
\item $D(u\|q)=\frac1K\sum_k\log\frac{1}{Kq_k}$. Each term with $k\ne y$ tends to $+\infty$ as $q_k\to 0$, so $D(u\|q)\to\infty$.

$D(q\|u)=\log K-H(q)\to\log K-0=\log K$, which is finite.

Hence $L_{\rm LS}(q)\to\infty$, since $(1-\varepsilon)(-\log q_y)\to 0$ but $\varepsilon D(u\|q)\to\infty$; and $L_{\rm CP}(q)\to 0+\beta\log K-\beta\log K=0$. Full confidence is infinitely expensive under label smoothing, but has finite loss under the confidence penalty.

%% (c) sol
\item Use the conventions $0\log\frac{0}{a}=0$ for $a\ge 0$ and $b\log\frac{b}{0}=+\infty$ for $b>0$.
\begin{enumerate}[i)]
\item $D(p\|q)=\sum_k p_k\log\frac{p_k}{q_k}$ is a finite sum. Terms with $p_k=0$ contribute $0$; a term with $p_k>0$ is finite if $q_k>0$ and $+\infty$ if $q_k=0$. Hence $D(p\|q)<\infty$ iff {\bf $q_k>0$ for every $k$ with $p_k>0$}.
\item $D(q\|p)=\sum_k q_k\log\frac{q_k}{p_k}$. Terms with $q_k=0$ contribute $0$; a term with $q_k>0$ is finite if $p_k>0$ and $+\infty$ if $p_k=0$. Hence $D(q\|p)<\infty$ iff {\bf $q_k=0$ for every $k$ with $p_k=0$}.
\item Since $u_k=\frac1K>0$ for every $k$:
\begin{itemize}
\item by (i), $D(u\|q)<\infty$ iff {\bf $q_k>0$ for every $k$} (and $D(u\|q)\to\infty$ as any $q_k\to 0$);
\item by (ii), the condition is empty, so $D(q\|u)<\infty$ for {\bf every} $q$; indeed $0\le D(q\|u)=\log K-H(q)\le\log K$.
\end{itemize}
\end{enumerate}
(Remark: this explains (b). $L_{\rm LS}$ contains $\varepsilon D(u\|q)$, which puts an unbounded cost on any class probability approaching $0$, so $L_{\rm LS}\to\infty$ as $q_y\to 1$. $L_{\rm CP}$ contains $\beta D(q\|u)$, which is bounded, so the fully confident limit has finite loss. Neither is a hard constraint: e.g.\ $L_{\rm LS}$ is minimized at $q=o^{(\varepsilon)}$, but any $q$ with all $q_k>0$ has finite loss. Minimizing $D(p\|q)$ over $q$ is called \emph{zero-avoiding} and minimizing $D(q\|p)$ \emph{zero-forcing}, after (i) and (ii).)
\end{enumerate}
\end{solution}

\begin{problem}
\newpage {\bf problem 2 cont.}
\end{problem}
\newpage


%% =====================================================================
%%  Problem 3
%% =====================================================================
\item {\bf Heavy Tails and the Median of Means} {\it (30 points)}

Let $X$ be a random variable with $X\ge 1$ and
\[
P(X>x)=x^{-3}\qquad\text{for all } x\ge 1.
\]
Let $X_1,\dots,X_n$ be i.i.d.\ copies of $X$, and $\bar X_n=\frac1n\sum_{i=1}^n X_i$. One can check that $\mathbb{E}[e^{tX}]=\infty$ for every $t>0$, so the Chernoff bound cannot be applied directly to $\bar X_n$.

{\it Hint.} For a nonnegative random variable $W$, $\mathbb{E}[W]=\int_0^\infty P(W>w)\,dw$. You may use this without proof.

\begin{enumerate}[(a)]
\item Show that $\mathbb{E}[X^2]=\int_0^\infty 2x\,P(X>x)\,dx$. Then compute $\mu=\mathbb{E}[X]$ and $\sigma^2={\rm Var}(X)$. {\it (5 points)}

\item For $a>0$, show that
\[
P(\bar X_n\ge\mu+a)\ \ge\ \big(n(\mu+a)\big)^{-3},
\]
and compare it with the upper bound given by Chebyshev's inequality. Conclude that $P(\bar X_n\ge\mu+a)$ decays polynomially, not exponentially, in $n$. [Hint: $\bar X_n\ge X_1/n$.] {\it (5 points)}

\item \emph{Median of means.} Split the $n$ samples into $k$ groups of $m$ samples each ($n=km$), compute the sample mean $\bar X^{(j)}$ of each group $j=1,\dots,k$, and let $\hat\mu$ be the median of the $k$ group means. Fix $a>0$ and call group $j$ \emph{bad} if $|\bar X^{(j)}-\mu|\ge a$. {\it (15 points)}
\begin{enumerate}[i)]
\item Show that if $m\ge 6/a^2$, each group is bad with probability at most $\frac18$. [Hint: Chebyshev's inequality for $\bar X^{(j)}$.] {\it (3 points)}
\item Show that if $|\hat\mu-\mu|\ge a$, then at least $k/2$ groups are bad. [Hint: if fewer than half of the groups are bad, where do more than half of the group means lie, and where is their median?] {\it (4 points)}
\item Let $Z$ be the number of bad groups. Since the groups use disjoint samples, you may assume the worst case $Z\sim{\rm Binomial}(k,\frac18)$. Apply the Chernoff bound to show that for every $t>0$,
\[
P(Z\ge k/2)\ \le\ \Big(e^{-t/2}\big(\tfrac78+\tfrac18 e^{t}\big)\Big)^k.
\]
[Hint: what is the moment generating function of a Bernoulli random variable?] {\it (4 points)}
\item Choose $e^t=7$ and show that the bound becomes $\big(\frac{\sqrt7}{4}\big)^k$. Combine (ii)--(iv) to conclude that with $m=\lceil 6/a^2\rceil$, $P(|\hat\mu-\mu|\ge a)$ decays exponentially in $n$. {\it (4 points)}
\end{enumerate}

\item Explain in two or three sentences why the median of means succeeds where the sample mean $\bar X_n$ fails. {\it (5 points)}
\end{enumerate}

%%% Solution of Problem 3
\begin{solution}
\vskip 0.2in
{\bf Solution:}

\begin{enumerate}[(a)]

%% (a) sol
\item Apply the hint to $W=X^2\ge 0$ and substitute $s=x^2$ ($ds=2x\,dx$):
\[
\mathbb{E}[X^2]=\int_0^\infty P(X^2>s)\,ds=\int_0^\infty P(X>\sqrt s)\,ds=\int_0^\infty 2x\,P(X>x)\,dx.
\]
Since $P(X>x)=1$ for $x<1$,
\[
\mu=1+\int_1^\infty x^{-3}\,dx=1+\tfrac12=\tfrac32,\qquad
\mathbb{E}[X^2]=1+\int_1^\infty 2x^{-2}\,dx=1+2=3,
\]
so $\sigma^2=3-\frac94=\frac34$. (Equivalently, use the density $f(x)=3x^{-4}$ for $x\ge 1$.)

%% (b) sol
\item All $X_i\ge 1>0$, so $\bar X_n\ge X_1/n$. If $X_1\ge n(\mu+a)$, then $\bar X_n\ge\mu+a$. Hence
\[
P(\bar X_n\ge\mu+a)\ \ge\ P\big(X_1\ge n(\mu+a)\big)=\big(n(\mu+a)\big)^{-3}.
\]
Chebyshev gives the upper bound
\[
P(\bar X_n\ge\mu+a)\le P(|\bar X_n-\mu|\ge a)\le\frac{\sigma^2}{na^2}=\frac{3}{4na^2}.
\]
So for fixed $a>0$ the probability tends to $0$ but is bounded above and below by polynomial rates ($n^{-1}$ and $n^{-3}$). In particular no bound $Ce^{-cn}$ with fixed $C,c>0$ can hold for all large $n$, because $(n(\mu+a))^{-3}\gg Ce^{-cn}$.

%% (c) sol
\item
\begin{enumerate}[i)]
\item $\bar X^{(j)}$ is the mean of $m$ i.i.d.\ copies of $X$, so ${\rm Var}(\bar X^{(j)})=\sigma^2/m$. By Chebyshev,
\[
P(|\bar X^{(j)}-\mu|\ge a)\le\frac{\sigma^2}{ma^2}=\frac{3}{4ma^2}\le\frac34\cdot\frac16=\frac18\qquad\text{when } m\ge 6/a^2.
\]
\item Suppose fewer than $k/2$ groups are bad. Then more than half of the group means lie in the interval $(\mu-a,\mu+a)$. If $\hat\mu\ge\mu+a$, then more than half of the group means would lie strictly below $\hat\mu$, contradicting the fact that a median has at least half of the values at or above it. The case $\hat\mu\le\mu-a$ is symmetric. Hence $|\hat\mu-\mu|<a$. Contrapositive: $|\hat\mu-\mu|\ge a$ implies at least $k/2$ bad groups, i.e., $\{|\hat\mu-\mu|\ge a\}\subseteq\{Z\ge k/2\}$.
\item Chernoff bound: for $t>0$, $e^{tz}$ is increasing, so by Markov's inequality
\[
P(Z\ge k/2)=P(e^{tZ}\ge e^{tk/2})\le e^{-tk/2}\,\mathbb{E}[e^{tZ}].
\]
$Z=\sum_{j=1}^k B_j$ with independent $B_j\sim\Bern(\frac18)$, and the MGF of a $\Bern(\frac18)$ variable is $\mathbb{E}[e^{tB_j}]=\frac78+\frac18e^t$. By independence, $\mathbb{E}[e^{tZ}]=\prod_j\mathbb{E}[e^{tB_j}]=\big(\frac78+\frac18e^t\big)^k$. Combining,
\[
P(Z\ge k/2)\le\Big(e^{-t/2}\big(\tfrac78+\tfrac18e^t\big)\Big)^k.
\]
\item With $e^t=7$: $e^{-t/2}\big(\frac78+\frac78\big)=\frac1{\sqrt7}\cdot\frac74=\frac{\sqrt7}{4}\approx 0.661<1$. By (ii) and (iii),
\[
P(|\hat\mu-\mu|\ge a)\le P(Z\ge k/2)\le\Big(\frac{\sqrt7}{4}\Big)^k.
\]
(If the bad-group probability is $q<\frac18$ rather than exactly $\frac18$, $Z$ is only smaller, so the bound still holds.) With $m=\lceil 6/a^2\rceil$ fixed, $k=n/m$ grows linearly in $n$, so the bound $\big(\frac{\sqrt7}{4}\big)^{n/m}$ decays exponentially in $n$. (The choice $e^t=7$ is the one that minimizes the bound in (iii).)
\end{enumerate}

%% (d) sol
\item A single sufficiently large sample can push $\bar X_n$ above $\mu+a$, and with a heavy tail such a sample appears with only polynomially small probability. The median of means confines any huge sample to a single group, where it spoils only one vote. The median fails only if half of the independent groups fail at once, which is a coin-flip large deviation with an exponentially small probability.
\end{enumerate}
\end{solution}

\begin{problem}
\newpage {\bf problem 3 cont.}
\end{problem}
\begin{problem}
\newpage {\bf extra sheet}
\newpage {\bf extra sheet}
\end{problem}

\end{enumerate}

\end{document}
